\documentclass[11pt]{article}
\usepackage[a4paper,margin=25mm]{geometry}
\usepackage[T1]{fontenc}
\usepackage{lmodern,microtype}
\usepackage{amsmath,amssymb,amsthm,mathtools,bm}
\usepackage{booktabs,array,enumitem}
\usepackage{xcolor}
\usepackage[colorlinks=true,linkcolor=blue!45!black,citecolor=blue!45!black,urlcolor=blue!45!black]{hyperref}
\usepackage{cleveref}
\hypersetup{
 pdftitle={Adjacency spectral likelihood classification in stochastic block models},
 pdfsubject={Constant oracle-risk bounds at every diverging expected-degree scale},
 pdfkeywords={stochastic block model, adjacency spectrum, Bernoulli likelihood, oracle risk}
}
\setlist{itemsep=3pt,topsep=5pt,leftmargin=2em}
\setlength{\emergencystretch}{2em}
\allowdisplaybreaks[2]
\newtheorem{theorem}{Theorem}[section]
\newtheorem{lemma}[theorem]{Lemma}
\newtheorem{proposition}[theorem]{Proposition}
\newtheorem{corollary}[theorem]{Corollary}
\theoremstyle{definition}
\newtheorem{definition}[theorem]{Definition}
\newtheorem{assumption}[theorem]{Assumption}
\theoremstyle{remark}
\newtheorem{remark}[theorem]{Remark}
\newcommand{\R}{\mathbb R}
\newcommand{\E}{\mathbb E}
\newcommand{\PP}{\mathbb P}
\newcommand{\one}{\mathbf 1}
\newcommand{\diag}{\operatorname{diag}}
\newcommand{\rank}{\operatorname{rank}}
\newcommand{\tr}{\operatorname{tr}}
\newcommand{\op}{\mathrm{op}}
\newcommand{\Ber}{\operatorname{Bern}}
\newcommand{\KL}{\operatorname{KL}}
\newcommand{\TV}{\operatorname{TV}}
\newcommand{\col}{\operatorname{col}}
\newcommand{\supp}{\operatorname{supp}}
\newcommand{\argmax}{\operatorname*{arg\,max}}
\newcommand{\argmin}{\operatorname*{arg\,min}}
\newcommand{\norm}[1]{\left\lVert#1\right\rVert}
\newcommand{\abs}[1]{\left\lvert#1\right\rvert}
\newcommand{\ip}[2]{\left\langle#1,#2\right\rangle}
\newcommand{\ind}[1]{\mathbf1\{#1\}}
\newcommand{\Mis}{\operatorname{Mis}}
\title{Adjacency spectral likelihood classification\\in stochastic block models}
\author{}
\date{21 October 2026}
\begin{document}
\maketitle
\begin{abstract}
We study community classification from the raw degrees and a single
eigensystem of a high-degree-regularized adjacency matrix. A square-root
cap reweights only vertices above a large empirical degree threshold
and permits arbitrarily slowly diverging expected degrees.
An exact degree correction, deterministic spectral initialization and
Bernoulli block likelihood updates define a finite classifier. For any
fixed number of communities, comparable community sizes and a positive
full-rank block profile, its expected misclassification proportion is at
most a constant times the oracle risk. The oracle knows the block
probabilities and all labels except that of the vertex being classified.
The bound holds whenever \(np\to\infty\) and \(np^2\) is bounded.
The proof combines truncated spectral-row energy, a local scalar
leave-one-out comparison and uniform fitted-score bounds. It imposes
no logarithmic lower bound on expected degree.
\end{abstract}
\section{Introduction}

The incident edges of a vertex determine a natural likelihood classifier
when the community parameters and the other labels are known. A spectral
method replaces those unknown labels by a low-dimensional representation
of the graph. We ask whether the degrees and a single adjacency
eigensystem retain enough information to match the error of the
incident-row likelihood classifier up to a constant factor.

We construct an explicit classifier with this property for stochastic
block models with arbitrary fixed \(K\). Its spectral input is a
regularized adjacency matrix \(\widetilde A=WAW\): the diagonal weight is one
unless the observed degree exceeds a sufficiently large fixed
multiple of the average degree. Above that threshold a square-root
cap controls the row energy. The decoder retains the raw degrees,
reconstructs a signed matrix with exactly those row sums, and alternates
symmetric Bernoulli block fitting with ordinary likelihood classification.
The only initialization is a deterministic radius-count procedure on
normalized spectral rows. At least one likelihood update is required;
any finite number of the specified updates is allowed.

When no degree exceeds the threshold, \(W=I\), so the encoding is
exactly the leading adjacency eigenpairs together with the raw degrees.
For the full sparse range the theorem uses the stated cap. The weighted
population signal \(W\Omega W\) has informative space \(\col(WZ)\), with one direction
per true community after row normalization. This explicit adjacency
construction is the version analyzed throughout the article.

The main result holds for every diverging degree scale \(d=np\), subject
to the upper sparsity condition \(np^2\le C_1\). In particular, it covers
\(d=\log\log n\) and slower divergence. The proof controls average
classification risk directly, allowing isolated vertices and exceptional
spectral rows. A graph-wide bound on the largest spectral row would
discard this range. Instead, a truncated-energy estimate bounds the
probability that the tested incident row actually encounters an
exceptional coefficient.

The target comparison is sharp at the level of risk. For a pair of
communities with exact Bernoulli Chernoff information \(I\), the oracle
error is of order \(d^{-1/2}e^{-I}\). Sharp Chernoff bounds in community
detection are developed by Zhou and Li~\cite{zhouli}; spectral error
exponents are studied by Zhang~\cite{zhang}. Here the local spectral
and fitted-likelihood errors are bounded on the scale needed to
preserve a constant risk multiplier. The geometric input uses sparse
graph regularization~\cite{llv2017} and, at larger degree scales,
a bounded-entry matrix norm inequality~\cite{bvh}.

\section{Model and the incident-row oracle}\label{sec:statements}

There are \(n\) vertices and a fixed number \(K\ge2\) of communities.
Let \(z\in[K]^n\), let \(Z\) be its indicator matrix, and let
\(n_a=\#\{i:z_i=a\}\). The symmetric adjacency matrix has zero diagonal
and independent upper-triangular entries satisfying
\begin{equation}\label{eq:mean-model}
 A_{ij}\sim\Ber(P_{z_i z_j}),\qquad
 P=pC,\qquad d=np,\qquad \Omega=ZPZ^\top,\qquad i<j.
\end{equation}
The structural constants obey
\begin{equation}\label{eq:structural}
 n_a\ge\pi_-n,\qquad C=C^\top,\qquad
 0<c_-\le C_{ab}\le c_+<\infty,\qquad
 \sigma_{\min}(C)\ge\kappa>0.
\end{equation}
The number \(K\), \(\pi_-\), \(c_-\), \(c_+\) and \(\kappa\) are fixed.
The matrix \(C\) may vary with \(n\) inside this class. Constants in
bounds may depend on these fixed quantities and on the stated upper
bound for \(np^2\), but not on \(n\). The completed mean \(\Omega\)
includes a diagonal solely for algebra; the actual mean is
\(\Omega-\diag(\Omega_{ii})\). We keep this \(O(p)\) difference in the
proof.

For a partition \(\widehat z\), define
\[
 \ell(\widehat z,z)=\frac1n\min_{\sigma\in\mathfrak S_K}
       \sum_i\ind{\sigma(\widehat z_i)\ne z_i}.
\]
The oracle knows \(P\) and all labels except the tested label. At vertex
\(i\) it maximizes the exact Bernoulli log likelihood
\begin{equation}\label{eq:sbm-oracle}
 L_{i,*}(a)=\sum_{j\ne i}
 \{A_{ij}\log P_{az_j}+(1-A_{ij})\log(1-P_{az_j})\}.
\end{equation}
It uses uniform candidate weights and a fixed rule for ties. An exact
global community-size constraint is not used to infer the missing
label. If \(z_{i,*}\) denotes this local classifier, its average
directional error is
\begin{equation}\label{eq:oracle-risk}
 R_{\mathrm{SBM},*}(n)=\frac1n\sum_i\PP\{z_{i,*}\ne z_i\}.
\end{equation}
This is the oracle used throughout the article.

\section{A decoder from degrees and adjacency spectrum}
\label{restore:algorithm}

The input is \(A\) and \(K\). Fix a sufficiently large constant
\(b_0>0\) and a sufficiently small constant \(c_0>0\), depending
only on the structural class. All choices and
ties below are deterministic. A fixed fallback partition defines the
algorithm whenever a required operation is undefined.

\paragraph{Encoding.}
Write
\[
 \mathbf d=A\one,\qquad D_i=\mathbf d_i,\qquad
 \bar D=\frac1n\sum_iD_i.
\]
\[
 \tau=b_0\bar D,\qquad G=\diag\{\max(D_i,\tau)\},\qquad
 W=\tau^{1/2}G^{-1/2},\qquad \widetilde A=WAW.
\]
The all-zero graph has the fixed output. Otherwise \(G\) is positive,
and the diagonal weights are
\[
 W_i=\begin{cases}
  1,&D_i\le\tau,\\
  \sqrt{\tau/D_i},&D_i>\tau.
 \end{cases}
\]
Thus edges between vertices below the threshold keep their original
weights. Retain the raw degrees and the \(K\) eigenpairs of \(\widetilde A\)
with largest
absolute eigenvalues:
\[
 U^\top U=I_K,\qquad \widetilde A U=U\Xi.
\]
Only \((\mathbf d,U,\Xi)\) are used after this step; \(\tau,G,W\)
are recovered from \(\mathbf d\). All informative eigenvalues are
retained by magnitude, allowing either sign.

\paragraph{Exact-degree reconstruction.}
Put
\[
 H=WUU^\top W^{-1},\qquad
 R_0=W^{-1}U\Xi U^\top W^{-1}=AH.
\]
If the cap is inactive, these reduce to \(H=UU^\top\) and
\(R_0=A_K\), the rank-\(K\) adjacency spectral reconstruction.
Let \(v=\mathbf d-R_0\one\), let \(R_1=R_0+v\one^\top/n\),
and let \(r\) be the vector of diagonal entries of \(R_1\). Define
\begin{equation}\label{restore:reconstruction}
 R^h=R_1+\frac{r\one^\top}{n-1}
              -\frac{n}{n-1}\diag(r).
\end{equation}
Then
\[
 R^h\one=\mathbf d,\qquad \diag(R^h)=0.
\]
The matrix \(R^h\) can be signed and need not be symmetric. Its entries
are not clipped. The formula is computed entirely from the retained
quantities. In particular, an isolated vertex has a zero reconstructed
row and does not stop classification of other vertices.

\paragraph{Initialization.}
Normalize each nonzero row of \(U\) to unit length; assign a fixed unit
direction to a zero row. Among the remaining rows, select one having
the largest number of remaining rows within distance \(1/8\), and
remove all rows within distance \(1/4\) of it. Repeat for \(K\) centers,
then assign every row to its nearest center. If fewer than \(K\)
centers can be selected, use the fixed fallback. This is a finite
procedure using \(O(Kn^2)\) distance comparisons.

\paragraph{Block fitting and classification.}
Given current labels \(\widehat z\), for \(a\le b\) let
\(\mathcal I_{ab}(\widehat z)\) contain all ordered pairs \(i\ne j\)
whose unordered label pair is \(\{a,b\}\). Let
\[
 E_{ab}=\sum_{(i,j)\in\mathcal I_{ab}(\widehat z)}R^h_{ij},
 \qquad N_{ab}=|\mathcal I_{ab}(\widehat z)|,
 \qquad \varepsilon=c_0\bar D/n.
\]
The symmetric fitted matrix is
\begin{equation}\label{eq:block-fit}
 \widehat P_{ab}=\widehat P_{ba}
   =\operatorname{clip}_{[\varepsilon,\,1/2]}(E_{ab}/N_{ab}).
\end{equation}
This maximizes the ordered-pair Bernoulli working criterion with a
symmetric block probability. The fit is interior on the event used in
the proof. Update all labels simultaneously by maximizing
\begin{equation}\label{eq:fitted-row-score}
 \widehat L_i(a)=\sum_{j\ne i}
 \{R^h_{ij}\log\widehat P_{a\widehat z_j}
   +(1-R^h_{ij})\log(1-\widehat P_{a\widehat z_j})\}.
\end{equation}
In particular, the full nonedge term is retained. There is no target
degree parameter to profile in this ordinary SBM algorithm.
Write \(\widehat L_i^{ab}=\widehat L_i(a)-\widehat L_i(b)\)
for a fitted pairwise contrast.

Perform at least one block fit and likelihood update. The output may
follow any finite number of these hard refits, including finite
data-dependent stopping. Empty fitted groups or undefined operations
use the stated fallback; the proof shows they do not occur on the
aggregate event. No leave-one-out graph, true label or true model
parameter is an algorithm input.

\section{Oracle-risk guarantee}\label{sec:main-result}

\begin{theorem}\label{restore:sbm-theorem}
Under \eqref{eq:mean-model}--\eqref{eq:structural}, suppose
\[
 d=np\longrightarrow\infty,\qquad np^2\le C_1<\infty.
\]
For sufficiently large fixed \(b_0\) and sufficiently small fixed
\(c_0>0\), the algorithm in
\cref{restore:algorithm}, after any finite positive number of the
specified likelihood updates, satisfies
\begin{equation}\label{eq:target}
 \E\ell(\widehat z,z)\le C R_{\mathrm{SBM},*}(n)
\end{equation}
for all sufficiently large \(n\). The same bound holds for finite
data-dependent stopping. The constant \(C\) is independent of \(n\)
and of the number of refits.
\end{theorem}

There is no condition of the form \(np\ge c\log n\).
The assumptions allow general symmetric positive block profiles and
comparable, unequal community sizes. They do not require a planted
partition or equal expected degrees. The condition \(np^2\le C_1\)
is retained as an upper sparsity assumption.

For the exact pairwise Bernoulli information \(I_{i,ab}\) defined in
\eqref{eq:sbm-information}, \cref{lem:sharp} gives
\[
 R_{\mathrm{SBM},*}(n)\asymp
 \frac{1}{n\sqrt d}\sum_i\sum_{b\ne z_i}e^{-I_{i,z_i b}}.
\]
Thus \eqref{eq:target} preserves the sharp prefactor as well as the
exponent. It asserts a bounded risk ratio; it does not assert that
the ratio tends to one. Since a nonzero loss is at least \(1/n\),
the theorem also implies exact recovery whenever
\(nR_{\mathrm{SBM},*}(n)\to0\).

\subsection{Proof structure}

The proof has four ingredients. First, regularized adjacency geometry
gives an informative spectral gap of order \(d\),
\(O(d^{-1/2})\) subspace error and \(O(n)\) weighted cap energy
\(\sum_i(G_i-\tau)^2/G_i\). The threshold lies above all population
degrees on the aggregate event, so this energy is bounded by the
weighted degree fluctuation. These aggregate bounds suffice for
the deterministic initialization and an invariant \(O(n/d)\)
misclassification basin.

Second, a same-threshold truncated-energy recursion controls large
spectral rows in expectation. By taking a fixed row threshold large
enough, the average probability of an exceptional tested row, or of
an incident edge hitting an exceptional training row, is at most
\(C_He^{-Hd}\) for any prescribed fixed \(H\). This estimate replaces
a simultaneous row bound.

Third, on the local regular event, the exact block fit has error
\(O(1/n)\), and every fitted row contrast differs by \(O(1)\) from its
true block contrast. A scalar leave-one-out identity replaces the
full-graph spectral action by a function of independent training
data and the tested incident row, again with \(O(1)\) error. Under
the oracle Chernoff tilt, the resulting coefficient error has bounded
mean and quadratic mass.

Finally, a sharp Bernoulli transfer preserves
\(d^{-1/2}e^{-I_{i,ab}}\). Exceptional coefficients are paid for only
when the incident row hits them. Polynomially accurate regularized
geometry suffices below an internal logarithmic proof threshold;
an exponentially accurate norm event supplies the complementary
range. These estimates together cover every \(d\to\infty\).
The appendices give the complete argument and the proof of
\cref{restore:sbm-theorem}.

\paragraph{Organization of related material.}
The companion notes contain the degree-normalized formulation, degree
correction, alternative earlier decoders, counterexamples and
supplementary derivations. None of
those notes is required for this article's theorem or proof.


\appendix
\section{A sharp Bernoulli likelihood transfer}\label{sec:sharp}

We give the probability step explicitly to distinguish a risk factor from
an exponential-rate factor. Sharp Chernoff comparisons in community
detection are also developed by \citet{zhouli}.
The argument below allows a perturbation measurable from independent
training data.

\begin{lemma}[Bounded tilted error preserves the sharp prefactor]
\label{lem:sharp}
Let $X_j$ be independent Bernoulli variables with probabilities $p_j$.
Let $T$ be an affine function of $(X_j)_j$, with linear coefficients
$w_j$, and put $\psi(\alpha)=\log\E e^{-\alpha T}$.
Suppose $\abs{w_j}\le C$, the minimizer $\alpha^*$ satisfies
$c\le \alpha^*\le1-c$ for a fixed $c>0$, and under the probability
measure $\mathbb{Q}$ proportional to $e^{-\alpha^*T}$,
\begin{equation}\label{eq:sharp-premises}
 \E_{\mathbb{Q}} T=0,\quad c{np}\le\operatorname{Var}_{\mathbb{Q}}(T)\le C{np},
 \quad\sum_jr_j\le C{np},
\end{equation}
where $r_j=\mathbb{Q}(X_j=1)$. Define
\[
 I=-\psi(\alpha^*)+\frac12\log(np).
\]
Within this lemma, $I$ is the exact transform-based exponent of $T$.
The added logarithm incorporates the square-root prefactor; hence
$I-\tfrac12\log(np)$ is the negative log transform at the optimizing
tilt. Its comparison with the full-size $I_{k\ell}$ of
Section~\ref{sec:oracle-scale} is made in the application below.
For every fixed $L<\infty$ and all sufficiently large ${np}$, uniformly over
$|x|\le L$,
\begin{equation}\label{eq:sharp-oracle}
 e^{-I}\lesssim\PP(T<x)
 \le\PP(T\le x)\lesssim e^{-I}.
\end{equation}
The implicit constants may depend on $L$ and the fixed constants in
the oracle hypotheses, but not on $n$ or on $x$ within this interval.

Let $\mathcal W$ be training data independent of $(X_j)_j$ and let
\[
 \widehat T=T+\varepsilon_0+\sum_j\varepsilon_jX_j
\]
have coefficients measurable with respect to $\mathcal W$. On a measurable
training event suppose
\begin{equation}\label{eq:tilt-error}
 \max_j|\varepsilon_j|\le R,\qquad
 \left|\varepsilon_0+\sum_jr_j\varepsilon_j\right|\le B_1,\qquad
 \sum_jr_j\varepsilon_j^2\le B_2.
\end{equation}
For fixed $R,B_2$ and every real $x$,
\begin{equation}\label{eq:sharp-transfer}
 \PP(\widehat T\le x\mid\mathcal W)
 \lesssim\exp\{\alpha^*x+B_1+e^R B_2/2\}\,e^{-I}.
\end{equation}
The implicit factor may depend on the fixed oracle bounds and on
$R,B_2$, but is uniform in $x,B_1,n$ and over the training event.
In particular $B_1$ can
be a nonnegative training-measurable random bound, provided its indicated
exponential moment is integrated when removing the conditioning.
\end{lemma}

\begin{proof}
Under $\mathbb{Q}$, the coordinates remain independent Bernoulli variables,
\[
 r_j=\frac{p_je^{-\alpha^*w_j}}{1-p_j+p_je^{-\alpha^*w_j}}.
\]
The sum of their centered weighted third absolute moments is at most
$C\sum_jr_j=O({np})$. Berry--Esseen and
\eqref{eq:sharp-premises} therefore show that every interval of length one
has $\mathbb{Q}$ probability at most $C/\sqrt{np}$, uniformly in
its location. The exact change of measure is
\[
 \frac{\mathrm{d}\PP}{\mathrm{d}\mathbb{Q}}
 =\exp\{\psi(\alpha^*)+\alpha^*T\}
 =\sqrt{np}\,e^{-I+\alpha^*T}.
\]
Partitioning $(-\infty,x]$ into unit intervals cancels the
$\sqrt{np}$ factor against the interval-probability bound and gives
the upper bound in \eqref{eq:sharp-oracle}.

For the lower bound choose a fixed interval $(x-L_0,x)$, with $L_0$
sufficiently large. Its Gaussian probability under a centered normal
variable of variance in $[c{np},C{np}]$ is at least $cL_0/\sqrt {np}$ for large
${np}$, uniformly over bounded $x$. Choose $L_0$ so this exceeds twice the
Berry--Esseen error. The $\mathbb{Q}$ probability is then at least $c'/\sqrt{np}$.
On this interval $e^{\alpha^*T}$ has a fixed positive lower bound,
so the same change-of-measure identity gives $c_Le^{-I}$.
Using strict endpoints also handles arbitrary lattice atoms and tie rules.

For the perturbation, independence of training and the row allows us to
regard $\varepsilon_0,\varepsilon$ as fixed. The elementary Bernoulli exponential-moment bound
gives
\begin{align*}
 \log\E e^{-\alpha^*\widehat T}
 &= -I+\tfrac12\log(np)-\alpha^*\varepsilon_0+\log\E_{\mathbb{Q}}e^{-\alpha^*\sum_j\varepsilon_jX_j}\\
 &\le -I+\tfrac12\log(np)-\alpha^*\left(\varepsilon_0+\sum_jr_j\varepsilon_j\right)
       +\frac{(\alpha^*)^2e^{\alpha^*R}}2\sum_jr_j\varepsilon_j^2\\
 &\le -I+\tfrac12\log(np)+B_1+e^R B_2/2.
\end{align*}
Now tilt the original law by $e^{-\alpha^*\widehat T}$.
Write $\widehat r_j=\widehat{\mathbb Q}(X_j=1)$.
Under this new measure $\widehat{\mathbb{Q}}$ the coordinates remain independent and
their odds are those under $\mathbb{Q}$ multiplied by $e^{-\alpha^*\varepsilon_j}$.
Consequently their probabilities and variances are within fixed factors
of those under $\mathbb{Q}$. Furthermore,
\[
 \left(\sum_jr_j(1-r_j)(w_j+\varepsilon_j)^2\right)^{1/2}
 \ge \sqrt{c{np}}-\sqrt{B_2}.
\]
For large ${np}$, the variance of $\widehat T$ under $\widehat{\mathbb{Q}}$ is therefore
between positive constant multiples of ${np}$. Its centered third absolute
moments sum to $O({np})$ because $w_j+\varepsilon_j$ are uniformly bounded and
$\sum_j\widehat r_j=O({np})$. Berry--Esseen bounds every unit interval under
$\widehat{\mathbb{Q}}$ by $C/\sqrt {np}$, irrespective of its mean. Multiplication by the displayed exponential-moment bound
cancels the $1/\sqrt{np}$ interval factor. The same interval summation
therefore proves \eqref{eq:sharp-transfer}.
\end{proof}

\subsection{Application to the actual Bernoulli SBM oracle}

Fix a tested vertex with true label $k$ and realize the count vector
from Section~\ref{sec:statements} using its actual incident edges.
For this fixed target, $T_{k\ell}$ is therefore exactly its raw-row
contrast, not an independent copy. Its distribution depends only on
$k,\ell$, the community sizes and $P$. The vertex is still indexed by
$i$ wherever its deleted graph or projection is used.
For the projection arguments the linear coefficient vector is
\begin{equation}\label{eq:sbm-weights}
 (w_{k\ell})_j
   =\log\frac{P_{kz_j}(1-P_{\ell z_j})}
                    {P_{\ell z_j}(1-P_{kz_j})},\qquad j\in[n].
\end{equation}
Thus, for the chosen target,
\[
 T_{k\ell}=\sum_{j\ne i}A_{ij}(w_{k\ell})_j
       +\sum_{j\ne i}\log\frac{1-P_{kz_j}}{1-P_{\ell z_j}}.
\]
The nonedge term is retained, and $w_{k\ell,-i}$ means that coordinate
$i$ is deleted. The independent binomial counts give the exact transform
\begin{equation}\label{eq:sbm-information}
 \log\E_k e^{-\alpha T_{k\ell}}
 =\sum_{r=1}^K(n_{r}-\ind{r=k})\,
 \log\!\left[P_{kr}^{1-\alpha}P_{\ell r}^\alpha
 +(1-P_{kr})^{1-\alpha}(1-P_{\ell r})^\alpha\right].
\end{equation}
The self-loop contributes nothing. To compare this exact transform
with the full-size expression in Section~\ref{sec:oracle-scale}, note
that the additional affinity factor obeys, for $0\le\alpha\le1$,
\[
 1-p\le P_{kk}^{1-\alpha}P_{\ell k}^{\alpha}
 +(1-P_{kk})^{1-\alpha}(1-P_{\ell k})^{\alpha}\le1.
\]
The lower bound follows from the nonedge term, and the upper bound
from H\"older's inequality. Taking negative logarithms and maximizing
shows that
\[
 0\le I_{k\ell}
       +\min_{0\le\alpha\le1}\log\E_k e^{-\alpha T_{k\ell}}
       -\tfrac12\log(np)\le-\log(1-p).
\]
Thus the exact transform-based exponent in \cref{lem:sharp} is
$I_{k\ell}+O(p)$, and its exponential lies between $e^{-I_{k\ell}}$
and $(1-p)^{-1}e^{-I_{k\ell}}$. Under $np^2\le C_1$, eventually
$p\le1/2$, so this additional multiplier is uniformly at most two.
Every optimizing tilt and change-of-measure identity in the lemma
still uses the exact transform \eqref{eq:sbm-information}, not a
fictitious observed self-loop. The optimizer of the full-size
expression need not coincide with that exact optimizer.
On the compact positive full-rank classes used here, probabilities are
comparable to $p$ and bounded away from one, the weights are bounded,
$\alpha^*$ stays away from the endpoints, and the tilted variance is comparable
to ${np}$. To verify the last assertions, the derivatives of the transform
at its endpoints are respectively minus and plus the two directional
Bernoulli KL divergences, both of order ${np}$. Its second derivative is
bounded by $C{np}$. The row separation implied by
$\sigma_K(P)\ge c p$ gives a lower bound $c{np}$ on that derivative
throughout $[0,1]$: tilted probabilities are comparable to $p$ and the
sum of squared log-odds differences is at least $cn$.
Thus \cref{lem:sharp} applies uniformly to every pair. Combining it
with the one-coordinate comparison above gives, for every fixed
bounded threshold $x$, the bounds
\[
 \PP_k(T_{k\ell}\le x)\asymp e^{-I_{k\ell}},
 \qquad -\log\PP_k(T_{k\ell}\le x)=I_{k\ell}+O(1).
\]
These are constant-factor comparisons, not exact risk identities.
The use of $\tfrac12\log(np)$ is justified by the tilted variance
being comparable to $np$ and the optimizer staying away from the
endpoints on the stated parameter class. It is not a universal
prefactor formula for vanishing signal separation.

\subsection{Fixed-\texorpdfstring{$K$}{K} comparison and label alignment}

A strict pairwise defeat of the true community forces an oracle
error; an error entails a non-strict defeat by some competitor.
Under the joint community-$k$ count law this gives
\[
 \begin{split}
 \max_{\ell\ne k}\PP_k(T_{k\ell}<0)
 &\le\PP_k\!\left\{
       \min_{\ell<k}T_{k\ell}\le0
       \ \text{or}\ \min_{\ell>k}T_{k\ell}<0\right\}\\
 &\le\sum_{\ell\ne k}\PP_k(T_{k\ell}\le0).
 \end{split}
\]
The sharp strict and non-strict tails, together with fixed $K$, imply
\begin{equation}\label{eq:oracle-pairs}
 R_{\mathrm{oracle}}
 \asymp\sum_{k=1}^K\frac{n_k}{n}
                    \sum_{\ell\ne k}e^{-I_{k\ell}}.
\end{equation}
The weighting $n_k/n$ is the original average over vertices, regrouped
by their true community. This is not a redefinition as a worst-pair
oracle. The argument requires no independence among competing tests.
Fitted partitions are aligned with the true groups by a single common
permutation on the good event. Uniform score comparisons under that
alignment imply
\[
 \E[\Mis(\widehat z,z)\ind{\mathcal G}]
 \le\frac1n\sum_i\PP(\text{aligned fitted error at }i,\mathcal G).
\]
The ensuing pairwise bounds apply to fixed true coefficient vectors; a
data-dependent matching does not become an independence assumption.

\subsection{Exceptional coefficients}

\begin{lemma}[Paying exceptional coefficients only when the row hits them]
\label{lem:sparse-truncated-transfer}
Assume the oracle hypotheses of \cref{lem:sharp}, and write $p_j$
for the un-tilted row probabilities. Suppose $p_j\le Cr_j$, as
holds here for bounded oracle likelihood coefficients. On an independent
training event $\mathcal E$, suppose
\[
 \left|\varepsilon_0+\sum_jr_j\varepsilon_j\right|\le B_1,\qquad \sum_jr_j\varepsilon_j^2\le B_2.
\]
Here $B_1,B_2$ are fixed constants. For any fixed $R>0$ and bounded
$x$, the exceptional coordinates can be charged directly:
\begin{equation}\label{eq:sparse-truncated-transfer}
 \begin{split}
 \PP(\widehat T\le x)
 \le{}& C_Re^{-I}+\PP(\mathcal E^c)\\
 &+\E\!\left[\ind{\mathcal E}
             \sum_{j:|\varepsilon_j|>R}p_j\right].
 \end{split}
\end{equation}
The last term is the expected probability bound for the tested row
hitting an exceptional coefficient. Since
\[
 \sum_{j:|\varepsilon_j|>R}p_j
 \le\frac{C}{R^2}
      \sum_jr_j\varepsilon_j^2\ind{|\varepsilon_j|>R},
\]
a corresponding bound for the truncated quadratic error also suffices.
Either expected term of order $e^{-I}$ is enough.
If $I\le Cnp$ for a fixed $C$, an exceptional-term bound
$C_se^{-snp}$ is sufficient for any fixed $s>C$.
\end{lemma}

\begin{proof}
For analysis only, set
$\varepsilon'_j=\varepsilon_j\ind{|\varepsilon_j|\le R}$.
Then
\[
 \left|\varepsilon_0+\sum_jr_j\varepsilon'_j\right|
 \le B_1+\frac1R\sum_{j:|\varepsilon_j|>R}r_j\varepsilon_j^2
 \le B_1+B_2/R,\qquad
 \sum_jr_j(\varepsilon'_j)^2\le B_2.
\]
The bounded-coefficient lemma bounds the truncated score's error by
$C_Re^{-I}$ on $\mathcal E$.
The two random scores agree whenever $X_j=0$ on every exceptional
coordinate. Conditional on training, the probability they differ is
at most $\sum_{j:|\varepsilon_j|>R}p_j$. Integrating this bound and
paying $\PP(\mathcal E^c)$ proves the result. The truncation is a
proof device, not an additional operation in the decoder.
\end{proof}


\section{Aggregate geometry for the capped adjacency matrix}
\label{sec:sparse-extension}

Let $Z_{ik}=\ind{z_i=k}$ be the community indicator matrix.
The completed matrix $ZPZ^\top$ includes its diagonal, whereas
$\E A=ZPZ^\top-\diag(P_{z_i z_i})$. Their difference has norm at
most $p$ and is retained in the bounds below. The completed row
sums $\sum_jP_{z_i z_j}$ are block constant and comparable to $np$.
In this appendix put $\tau=b_0\bar D$ and
$G=\diag\{\max(D_i,\tau)\}$, with diagonal entries $G_i$.
Then $W=\tau^{1/2}G^{-1/2}$ is the algorithm's high-degree cap,
with $b_0$ sufficiently large for the structural class. Thus
$\widetilde A=WAW$ leaves every edge between vertices of degree at most $\tau$
unchanged. For analysis only, write
$Q=\widetilde A/\tau=G^{-1/2}AG^{-1/2}$ and
$\norm v_G^2=v^\top Gv$. The scalar rescaling $Q$ has the same
eigenvectors as the encoded adjacency matrix $\widetilde A$.

\begin{proposition}[Aggregate capped-adjacency geometry]\label{prop:sparse-floor}
Under \eqref{eq:structural}, let $np\to\infty$ and $np^2\le C_1$.
For every fixed $r>0$, with probability at least $1-C_rn^{-r}$,
$c{np}\le\tau\le C{np}$, $\tau\ge\max_i\sum_jP_{z_i z_j}$,
$c{n^2p}\le\sum_iD_i$, and
\begin{align}
 \norm{W(A-{ZPZ^\top})W}_{\op}&\le C_r\sqrt {np},
 &\sum_iG_i&\le C_r{n^2p},\label{eq:floor-noise}\\
 \sum_i\frac{(D_i-\sum_jP_{z_i z_j})^2}{G_i}
 +\sum_i\frac{(G_i-\tau)^2}{G_i}
 &\le C_rn.\label{eq:floor-degree-energy}
\end{align}
In particular,
$\norm{G^{-1/2}(A-{ZPZ^\top})G^{-1/2}}_{\op}\le C_r/\sqrt {np}$.
Let $U$ contain the $K$ eigenvectors of $\widetilde A$ with largest absolute
eigenvalues, and put $\Pi=UU^\top$ and $H=G^{-1/2}\Pi G^{1/2}$.
The eigenvalue equation gives $\widehat A=AH$ for the reconstruction
$\widehat A=W^{-1}U\Lambda U^\top W^{-1}$ in \cref{restore:algorithm}.
The retained eigenvalues of $\widetilde A$ have magnitude at least $c_r{np}$;
the remaining eigenvalues have magnitude at most $C_r\sqrt {np}$.
There is an orthogonal matrix $O$ such that the orthonormal population
basis
\[
 V=G^{-1/2}Z(Z^\top G^{-1}Z)^{-1/2}O
\]
satisfies
\[
 \norm{U-V}_F\le C_r/\sqrt {np},
 \qquad \sqrt{G_i}\norm{V_{i\cdot}}_2\le C_r\sqrt{p}.
\]
Moreover,
\begin{equation}\label{eq:floor-block-residual}
 \norm{(I-\Pi)G^{1/2}Z}_F\le C_r\sqrt n,
\end{equation}
and every block-constant $f$ with $\norm f_\infty\le C_f$ obeys
\begin{equation}\label{eq:floor-weighted-errors}
 \norm{(I-H)f}_G\le C_rC_f\sqrt n,\qquad
 \norm{(I-H)f}_2^2\le \frac{C_rC_f^2}{p}.
\end{equation}
The same assertions hold simultaneously after deletion of any fixed
number of vertices, with each cap recomputed from the remaining graph,
after increasing the probability parameter. No raw minimum-degree
condition is imposed.
\end{proposition}

\begin{proof}
The Bernoulli transform of the total number of edges gives
$c{n^2p}\le\sum_iD_i\le C{n^2p}$ except with probability $e^{-c{n^2p}}$.
For large enough $b_0$, this event implies
\[
 c{np}\le\tau=b_0\bar D\le C{np},
 \qquad \tau\ge\max_i\sum_jP_{z_i z_j}.
\]
Since $\E A_{ij}\le p$, put
$W_{h,ii}=\sqrt{\min(1,2{np}/D_i)}$, with value one at zero degrees.
The changed set $\mathcal S=\{i:D_i>2{np}\}$ contains at most
$\frac{10}{p}$ vertices except with probability $e^{-cn}$. To check the
strength of this bound, set $m=\lfloor\frac{10}{p}\rfloor+1$. If $m\le n$
and $m$ vertices exceed the threshold, some set $T$ of size $m$ has
$\sum_{i\in T}D_i>2m{np}$. This is an independent-edge sum with
coefficients in $\{1,2\}$ and mean at most $m{np}$. At
$t=(\log2)/2$, its exponential transform bounds the tail by
\[
 \exp\{-(\log2-1/2)m{np}\},
\]
because $(e^{tc}-1)/c\le(e^{2t}-1)/2=1/2$ for $c=1,2$.
Since $\log\binom nm\le m\log(en/m)=o(m{np})$, a union over $T$
gives the stated bound. If $m>n$, the claim is automatic.

Every row of $W_hAW_h$ has squared Euclidean norm at most
$W_{h,ii}^2D_i\le2{np}$. The row-norm form of the regularized adjacency
theorem of \citet[Theorem 2.1 and Remark 2.2]{llv2017} therefore yields
\[
 \norm{W_hAW_h-\E A}_{\op}\le C_r\sqrt {np}
\]
with failure probability at most $C_rn^{-r}$. The theorem permits
these graph-dependent weights. Also
\[
 \norm{W_h{ZPZ^\top} W_h-{ZPZ^\top}}_{\op}
 \le Cp\sqrt{n|\mathcal S|}\le C\sqrt {np},
\]
by the Frobenius norm: the difference is supported on rows or columns
in $\mathcal S$. The completed-mean diagonal has norm $O(p)$.
Consequently $\norm{W_h(A-{ZPZ^\top})W_h}_{\op}\le C_r\sqrt {np}$.
The exact diagonal comparison is
\[
 G^{-1/2}=(2{np})^{-1/2}
 \diag\!\left\{\sqrt{\frac{\max(D_i,2{np})}{G_i}}\right\}W_h.
\]
The displayed diagonal multiplier is bounded, since
$G_i\ge\max(D_i,c{np})$.
Conjugation gives the stated bound for the noise in $Q$;
multiplication by $\tau\asymp {np}$ proves the adjacency-scale bound
in \eqref{eq:floor-noise}. Since $G_i\le D_i+\tau$, the total
weight $\sum_iG_i$ is at most $C{n^2p}$.

Applying the scalar-rescaled noise to $G^{1/2}\one$ gives
\[
 \norm{G^{-1/2}(A-{ZPZ^\top})\one}_2
 \le C_r{(np)}^{-1/2}\norm{G^{1/2}\one}_2\le C_r\sqrt n,
\]
which is the first sum in \eqref{eq:floor-degree-energy}.
For the second sum, $G_i-\tau=0$ when $D_i\le\tau$; otherwise
$0\le G_i-\tau=D_i-\tau\le D_i-\sum_jP_{z_i z_j}$. Thus, coordinatewise,
\[
 \frac{(G_i-\tau)^2}{G_i}
 \le\frac{(D_i-\sum_jP_{z_i z_j})^2}{G_i}.
\]
The same bound also gives
$\sum_i(1-W_{ii})^2\le \frac{C_r}{p}$, since
$(1-W_{ii})^2\le(G_i-\tau)^2/(\tau G_i)$.

Let $J=Z^\top G^{-1}Z$. For every block $k$, Cauchy--Schwarz,
comparable group sizes, $G_i\ge c{np}$ and the total-degree bound give
\[
 \frac{c}{p}\le
 \frac{n_k^2}{\sum_{i:z_i=k}G_i}
 \le J_{kk}=\sum_{i:z_i=k}G_i^{-1}\le\frac{C}{p}.
\]
The nonzero spectrum of the capped population matrix $W{ZPZ^\top} W$ is
the spectrum of $\tau J^{1/2}PJ^{1/2}$. The bound
$\sigma_K(P)\ge c p$ gives a separation of order ${np}$ from zero.
The adjacency noise bound and
subspace perturbation then give the stated eigenspace and eigenvalue
bounds. The formula for $V$
and the lower bound on $J$ give its weighted row bound.

The matrix $\tau G^{-1/2}Z$ belongs to the population signal space
$\col(WZ)=\col(G^{-1/2}Z)$. Furthermore,
\[
 \norm{G^{1/2}Z-\tau G^{-1/2}Z}_F^2
 =\sum_i\frac{(G_i-\tau)^2}{G_i}\le C_rn,
 \qquad \norm{G^{1/2}Z}_F^2\le C_r{n^2p}.
\]
Together with the projector error $O({(np)}^{-1/2})$, these imply
\eqref{eq:floor-block-residual}. The matrix $H$ is the
$G$-orthogonal projector corresponding to $\Pi$, so the first bound
in \eqref{eq:floor-weighted-errors} follows by applying the residual bound to the $K$ block values of $f$.
The second uses $G_i\ge c{np}$.

Deleting a fixed number of vertices preserves the structural bounds.
Apply the same argument to each recomputed graph and choose a larger
fixed probability parameter before the finite-depth union bound.
\end{proof}

\section{Truncated spectral-row energy}
\label{restore:tail-section}

The next lemma controls large spectral rows in expectation. Its proof
uses a recurrence at a single truncation threshold, together with
concentration of independent deleted-row innovations.

Work under the fixed structural assumptions of
\cref{prop:sparse-floor}, with $np\to\infty$ and $np^2\le C_1$.
For the high-degree cap in \cref{restore:algorithm}, recall
\[
 \tau=b_0\bar D,\qquad
 G_j=\max(D_j,\tau),\qquad
 W=\tau^{1/2}G^{-1/2},\qquad \widetilde A=WAW.
\]
Let $U$ be an orthonormal basis of the informative $K$-dimensional
eigenspace and write $\widetilde AU=U\Lambda$.
All bounds below are uniform in the fixed compact structural class.
In particular,
\begin{equation}
 0\le{P_{z_i z_j}}\le Cp,\qquad
 \norm{{(ZPZ^\top)_{i\cdot}}}_2\le C{np}/\sqrt n.
 \label{restore:tail-population-rows}
\end{equation}

Fix a probability parameter $r>0$.  Apply
\cref{prop:sparse-floor} with a sufficiently larger fixed parameter to
the full graph and all graphs obtained by deleting at most two vertices.
Let $\mathcal E_r$ be the intersection of the resulting geometry and
two-sided total-degree events.  Restricting the population matrix
preserves the fixed structural bounds and comparable block masses.
Thus
\begin{equation}
 \PP(\mathcal E_r^c)\le C_rn^{-r}.
 \label{restore:tail-geometry-probability}
\end{equation}
On this event, for every graph in the deletion family, its recomputed
threshold is comparable to ${np}$, its total degree is at most $C_r{n^2p}$,
its centered weighted noise has norm at most $C_r\sqrt {np}$, its
informative eigenvalues have absolute values at least $c_r{np}$, and its
remaining eigenvalues have absolute values at most $C_r\sqrt {np}$.
For a deleted training graph, use the analogous event formed from
that graph and its further deletions. This event is measurable from
the training graph alone and has the same probability bound.

The following deterministic properties of the cap will be useful:
\begin{equation}
 W_j\le1,\qquad W_j^2D_j\le\tau,
 \qquad \norm{\widetilde A}_{\op}\le\tau.
 \label{restore:tail-floor-operator}
\end{equation}
Indeed, $G\ge\diag(A\one)$ and the inequality
$2|ab|\le a^2+b^2$ gives
$|x^\top Qx|\le\sum_j(D_j/G_j)x_j^2\le\norm x_2^2$,
where $Q=\widetilde A/\tau$.

\begin{lemma}[Truncated energy for the capped-adjacency eigenvectors]
\label{restore:tail-energy-lemma}
There are fixed constants $R_0,C,c>0$ such that, for every fixed
$R\ge R_0$ and all sufficiently large $n$,
\begin{equation}
 \E\!\left[\ind{\mathcal E_r}
       \sum_{j=1}^n\norm{U_{j\cdot}}_2^2
                  \ind{\sqrt n\norm{U_{j\cdot}}_2>R}\right]
 \le CR^2e^{-c{np}\log R}.
 \label{restore:tail-energy}
\end{equation}
The constants $R_0,C,c$ depend on the fixed structural bounds,
$b_0$, and $r$, but not on $R$ or $n$.  The sufficiently-large-$n$
condition may depend on the chosen fixed $R$.  The same result holds
for each recomputed deleted graph with its training-measurable
geometry event.
\end{lemma}

\begin{proof}
We give the weight perturbation, the same-threshold energy recurrence,
and its expectation bounds separately.

\medskip\noindent\textbf{Recomputed weights and the eigenspace residual.}
For $J=[n]\setminus\{i\}$, let $W^{(-i)}$ and $U^{(-i)}$ be
computed from $A_{JJ}$, with $U^{(-i)}$ embedded by a zero row at $i$.
Outside the geometry event choose an arbitrary deterministic
orthonormal fallback depending only on $A_{JJ}$.
For an all-zero training graph set $W^{(-i)}=I$ and use that
orthonormal fallback. Thus the proof-only completion always has
$\norm{W^{(-i)}}_{\op}\le1$, and the independent innovation below
is defined for every training graph.
Put $(\Gamma_i)_{jj}=W_j^{(-i)}/W_j-1$ for $j\in J$ and
$(\Gamma_i)_{ii}=0$.
The exact threshold difference is
\[
 \tau^{(-i)}-\tau
 =b_0\frac{\bar D-2D_i}{n-1}.
\]
Comparing the two square-root caps, using the comparability of the
thresholds to ${np}$, gives on $\mathcal E_r$
\begin{equation}
 |(\Gamma_i)_{jj}|
 \le C\left\{\frac{A_{ij}}{G_j}
             +\frac{{np}+D_i}{{n^2p}}\right\},
 \qquad \norm{\Gamma_i}_{\op}\le C/(np).
 \label{restore:tail-weight-change}
\end{equation}
This also covers a crossing of the cap threshold. It follows either
by splitting the cases in the two maxima or by comparing the logarithms
of $\sqrt{\tau/\max(D_j,\tau)}$; the relative threshold change is
$O(({np}+D_i)/({n^2p}))$.  The last bound uses $D_i\le n$ and ${np}\le Cn$.

There is an orthogonal matrix $O_i$ such that
\begin{equation}
 \norm{U-U^{(-i)}O_i}_F
 \le C\norm{U_{i\cdot}}_2+C\norm{\Gamma_iU}_F.
 \label{restore:tail-loo-residual}
\end{equation}
To verify this, write
$\widetilde A=\left(\begin{smallmatrix}0&m_i^\top\\m_i&B_i\end{smallmatrix}\right)$.
By \eqref{restore:tail-floor-operator},
$\norm{m_i}_2\le\sqrt\tau$ and $\norm{B_i}_{\op}\le\tau$.
Also $B_iU_J=U_J\Lambda-m_iU_{i\cdot}$.
The lower block of the recomputed deleted matrix is
$(I+\Gamma_i)B_i(I+\Gamma_i)$.
Subtracting the full matrix and applying the difference to $U$ gives
a Frobenius residual bounded by
$C{np}(\norm{U_{i\cdot}}_2+\norm{\Gamma_iU}_F)$.
The full informative spectrum is separated from the deleted
noninformative spectrum by $c{np}$.
The complementary-space Sylvester equation and orthogonal Procrustes
alignment give \eqref{restore:tail-loo-residual}.
Only separation from the complementary spectrum is used; the
informative eigenvalues can repeat or have either sign.

\medskip\noindent\textbf{Truncation and summable weight changes.}
Set $t=R/\sqrt n$.  Let $f_t$ be radial projection onto the Euclidean
ball of radius $t$, applied to each row, and put
\[
 \bar U=f_t(U),\qquad E_U=U-\bar U,\qquad
 \mathcal T_R(U)=\sum_j\norm{U_{j\cdot}}_2^2
                         \ind{\norm{U_{j\cdot}}_2>t}.
\]
The radial projection is $1$-Lipschitz and commutes with orthogonal
rotations.  Moreover $\norm{E_U}_F^2\le \mathcal T_R(U)$,
$\max_j\norm{\bar U_{j\cdot}}_2\le t$, and $\norm{\bar U}_F\le\sqrt K$.
Define
\[
 \xi_i=W_i\left(\sqrt{D_i}+C{np}/\sqrt n\right).
\]
By \eqref{restore:tail-floor-operator}, $\xi_i^2\le C{np}$.
Equation \eqref{restore:tail-weight-change} yields
\[
 \norm{\Gamma_i\bar U}_F
 \le Ct\frac{\sqrt{D_i}}{np}+C\frac{{np}+D_i}{{n^2p}},
\]
and therefore
\begin{equation}
 \frac{C\xi_i}{np}\norm{\Gamma_i\bar U}_F
 \le Ct\left\{\frac1{np}+\frac{\sqrt{D_i}}{{(np)}^{3/2}}\right\}.
 \label{restore:tail-bounded-weight}
\end{equation}
For the global term, use $t\ge n^{-1/2}$, $D_i\le n$, and
${np}\le Cn$: $D_i/(n{(np)}^{3/2})\le t\sqrt{D_i}/{(np)}^{3/2}$ and
$1/(n\sqrt {np})\le Ct/(np)$.
The square-root cap does not give $W_iD_i\le C{np}$. The
right side of \eqref{restore:tail-bounded-weight} can exceed a
sufficiently small fixed multiple of $t$, for large ${np}$, only if
$D_i\ge c_*{(np)}^3$.

The unbounded part satisfies the summable estimate
\begin{equation}
 \sum_i\xi_i^2\norm{\Gamma_iE_U}_F^2\le C\norm{E_U}_F^2.
 \label{restore:tail-summable-weight}
\end{equation}
Indeed, after interchanging $i,j$, the local coefficient of
$\norm{(E_U)_{j\cdot}}_2^2$ is at most
$C{np}D_j/G_j^2\le C$.
The global coefficient is bounded by
\[
 \frac{C}{n^2{np}}\sum_i({np}+D_i)^2\le C,
\]
because $\sum_iD_i\le C{n^2p}$ and
$\sum_iD_i^2\le n\sum_iD_i$.
These estimates do not use a maximum-degree event.

\medskip\noindent\textbf{Independent innovations and the recurrence.}
Define the unrotated, training-measurable-coefficient innovation
\begin{equation}
 K_i=(A_{iJ}-{(ZPZ^\top)_{iJ}})W^{(-i)}f_t(U_J^{(-i)}).
 \label{restore:tail-innovation}
\end{equation}
Let $E_0=W(A-{ZPZ^\top})W$.
The eigenvector equation is
\[
 U=W{ZPZ^\top} WU\Lambda^{-1}+E_0\bar U\Lambda^{-1}+E_0E_U\Lambda^{-1}.
\]
The first term has every row bounded by $C/\sqrt n$, by
\eqref{restore:tail-population-rows}.
For the second, compare $W_J\bar U_J$ to
$W^{(-i)}f_t(U_J^{(-i)})O_i$ and use radial Lipschitzness,
\eqref{restore:tail-loo-residual}, and
\eqref{restore:tail-weight-change}.  This gives
\begin{align}
 \norm{U_{i\cdot}}_2\le{}&
 C/\sqrt n+(C/(np))\norm{K_i}_2
 +(C\xi_i/(np))\norm{U_{i\cdot}}_2\nonumber\\
 &+(C\xi_i/(np))\norm{\Gamma_i\bar U}_F
 +(C\xi_i/(np))\norm{\Gamma_iE_U}_F
 +(C/(np))\norm{(E_0E_U)_{i\cdot}}_2.
 \label{restore:tail-row-recursion}
\end{align}
The completed-mean diagonal contributes at most $Cpt/(np)\le Ct/n$
and is absorbed in the bounded terms; it is not treated as an
observed self-loop.
Although $O_i$ may depend on the tested row, the argument only uses
$\norm{K_iO_i}_2=\norm{K_i}_2$.
No conditional concentration is applied to rotated coordinates.

Take $R$ sufficiently large and then ${np}$ sufficiently large to absorb
the baseline and self terms.  Retain from
\eqref{restore:tail-bounded-weight} the excess
$Ct\sqrt{D_i/{(np)}^3}\ind{D_i\ge c_*{(np)}^3}$.
For $\norm{U_{i\cdot}}_2>t$, separate the cases
$\norm{K_i}_2>c'{np}t$ and $\norm{K_i}_2\le c'{np}t$ and square.
The resulting bound is
\begin{align}
 \norm{U_{i\cdot}}_2^2\ind{\norm{U_{i\cdot}}_2>t}
 \le{}&\frac C{{(np)}^2}\norm{K_i}_2^2
                         \ind{\norm{K_i}_2>c'{np}t}\nonumber\\
 &+\frac C{{(np)}^2}\norm{(E_0E_U)_{i\cdot}}_2^2
 +\frac{C\xi_i^2}{{(np)}^2}\norm{\Gamma_iE_U}_F^2
 +\frac{Ct^2}{{(np)}^3}D_i\ind{D_i\ge c_*{(np)}^3}.
 \label{restore:tail-pointwise-recursion}
\end{align}
Summing, using $\norm{E_0}_{\op}^2\le C{np}$ and
\eqref{restore:tail-summable-weight}, and multiplying by
$\ind{\mathcal E_r}$ gives
\begin{align}
 \mathcal T_R(U)\ind{\mathcal E_r}
 \le{}&\frac C{{(np)}^2}\sum_i\norm{K_i}_2^2
                         \ind{\norm{K_i}_2>c'{np}t}
 +\frac C{np}\mathcal T_R(U)\ind{\mathcal E_r}\nonumber\\
 &+\frac{Ct^2}{{(np)}^3}\sum_iD_i\ind{D_i\ge c_*{(np)}^3}.
 \label{restore:tail-energy-recursion}
\end{align}
Both energy terms use the same threshold $R$.

\medskip\noindent\textbf{Expectation bounds.}
Conditional on $A_{JJ}$, \eqref{restore:tail-innovation} is a sum
of independent centered vectors.  In each coordinate its summands
are bounded by $t$, and its total variance is at most
\[
 Cp\norm{W^{(-i)}f_t(U^{(-i)})}_F^2\le Cp.
\]
These bounds hold for every training graph, without conditioning on
$\mathcal E_r$.  Bennett's inequality in the fixed $K$ coordinates,
followed by integration of the squared-norm tail, gives
\begin{equation}
 \E\!\left[\norm{K_i}_2^2
                        \ind{\norm{K_i}_2>c'{np}t}\right]
 \le C{(np)}^2t^2e^{-c_1{np}\log R}.
 \label{restore:tail-innovation-moment}
\end{equation}
To see the exponent, set $v=Cp$ and
$h(u)=(1+u)\log(1+u)-u$.
At $x=c'{np}t$, the Bennett exponent
$(v/t^2)h(tx/v)$ is at least $c_1{np}\log R$, since
$tx/v$ is proportional to $R^2$.
Beyond that point the derivative of the exponent is at least
$c\log R/t$, which gives the same bound upon tail integration.

Every $D_i$ is a Bernoulli sum with mean at most $C{np}$.
Its Chernoff bound and tail integration yield
\begin{equation}
 \E[D_i\ind{D_i\ge c_*{(np)}^3}]
 \le C{(np)}^3e^{-c_2{(np)}^3\log {np}}.
 \label{restore:tail-large-degree-moment}
\end{equation}
The left side is zero if $c_*{(np)}^3>n$.
Taking expectations in \eqref{restore:tail-energy-recursion},
absorbing $C/(np)$, and using $nt^2=R^2$ now gives
\[
 \E[\mathcal T_R(U)\ind{\mathcal E_r}]
 \le CR^2\{e^{-c_1{np}\log R}+e^{-c_2{(np)}^3\log {np}}\}.
\]
For each fixed $R$, the second exponential is eventually bounded by
the first after reducing $c_1$.  This proves
\eqref{restore:tail-energy}.  The proof is uniform under deletion
of a fixed number of vertices and uses their recomputed weights,
so it also proves the asserted training-graph version.
The same proof applies to any event supplying the stated aggregate
geometry, including the exponentially accurate event of
\cref{restore:high-aggregate}. The innovation and degree-tail bounds
are unconditional before restriction to that event.
\end{proof}

\begin{corollary}[Exceptional coefficients are hit with exponentially small mass]
\label{restore:tail-hit-corollary}
On a training graph with the preceding cap, put
$H=G^{-1/2}UU^\top G^{1/2}$.
Let $f$ be any training-measurable vector satisfying
$\norm f_\infty\le C_f$, and let $\varepsilon=(H-I)f$.
For every prescribed fixed $s_0>0$, there is a fixed $L_\varepsilon<\infty$
such that, for any independent tested Bernoulli row with probabilities
$p_j\le Cp$,
\begin{equation}
 \E\!\left[\ind{\mathcal E_r}
                 \sum_{j:|\varepsilon_j|>L_\varepsilon}p_j\right]
 \le C_{s_0}e^{-s_0{np}}.
 \label{restore:tail-hit-mass}
\end{equation}
Here $\mathcal E_r$ is the training-measurable geometry event.
\end{corollary}

\begin{proof}
On that event, $G_j\ge c{np}$ and $\sum_jG_j\le C{n^2p}$.
For a fixed constant $C$ determined by these two bounds,
\[
 |\varepsilon_j|
 \le C_f+\frac{\norm{U_{j\cdot}}_2}{\sqrt{G_j}}
                  \norm{U^\top G^{1/2}f}_2
 \le C_f+CC_f\sqrt n\norm{U_{j\cdot}}_2.
\]
For $L_\varepsilon=C_f+CC_fR$, the event
$|\varepsilon_j|>L_\varepsilon$ implies $\sqrt n\norm{U_{j\cdot}}_2>R$.
Thus \cref{restore:tail-energy-lemma} gives
\begin{equation}
 \E\!\left[\ind{\mathcal E_r}
                    \sum_j\ind{|\varepsilon_j|>L_\varepsilon}\right]
 \le\frac n{R^2}\E[\mathcal T_R(U)\ind{\mathcal E_r}]
 \le Cne^{-c{np}\log R}.
 \label{restore:tail-coefficient-count}
\end{equation}
Multiplication by $Cp$ bounds the expected hit mass by
$C{np} e^{-c{np}\log R}$.
Choose a fixed $R$ with $c\log R>s_0+1$ and absorb the factor ${np}$.
Independence is required only to interpret the conditional probability
that the tested row hits the training exceptional set; no independence
between the training geometry event and its eigenvectors is asserted.
\end{proof}

\begin{corollary}[Average exceptional target-row mass]
\label{restore:tail-target-corollary}
For every fixed degree cutoff $L_D$ and every fixed $s_0>0$, there is
a fixed $L_0$ such that
\begin{equation}
 \frac1n\sum_i\PP\!\left\{
    \sqrt{G_i}\norm{U_{i\cdot}}_2>L_0\sqrt{p},\quad
    D_i\le L_D{np},\quad \mathcal E_r\right\}
 \le C_{s_0,L_D}e^{-s_0{np}}.
 \label{restore:tail-target-mass}
\end{equation}
\end{corollary}

\begin{proof}
On the displayed event, $G_i\le C_{L_D}{np}$.
Choose $L_0=R\sqrt{C_{L_D}}$.
The target-row exception then implies
$\sqrt n\norm{U_{i\cdot}}_2>R$.
Apply \eqref{restore:tail-energy} and divide its corresponding
expected count by $n$.  Choosing fixed $R$ sufficiently large gives
\eqref{restore:tail-target-mass}.
\end{proof}


\section{Local stability under deletion of the tested vertex}
\label{restore:loo}

The comparison below requires spectral row bounds only on neighbors
of the tested vertex. The cap threshold is recomputed after deletion.
We divide each capped adjacency matrix by its own scalar threshold
to write a generalized eigenvalue equation; this does not change its
retained eigenspace or the reconstruction $AH$.

\begin{lemma}[A local scalar comparison for capped adjacency]\label{restore:loo-scalar}
Put the tested vertex $i$ first and write $A=\left(\begin{smallmatrix}0&x^\top\\x&B\end{smallmatrix}\right)$,
$D_i=\sum_jx_j\le L_D{np}$, and let
$\Delta=\diag\{\max(D_{B,j},\tau_B)\}$ be recomputed from $B$.
Suppose
\[
 c{np}\le\Delta_j,\qquad \sum_j\Delta_j\le C{n^2p},
 \qquad Q_B=\Delta^{-1/2}B\Delta^{-1/2}
   =U_B(\Lambda_B/\tau_B)U_B^\top+N,
\]
where $\Lambda_B$ contains the retained eigenvalues of $\tau_BQ_B$,
$NU_B=0$, $\norm N_{\op}\le C/\sqrt {np}$, and the $K$ retained
eigenvalues of $Q_B$ have a fixed gap from zero. Let the full diagonal cap weight be
$G=\diag(G_i,\Delta+E)$, with
\[
 c{np}\le G_i\le C{np},\qquad E=E_\mathrm{loc}+E_\mathrm{glob},\qquad
 0\le E_\mathrm{loc}\le\diag(x),\quad
 \norm{E_\mathrm{glob}}_{\op}\le Cp.
\]
Here ${np}\to\infty$ and ${np}\le Cn$. Suppose only on hit coordinates that
\begin{equation}\label{restore:loo-hit}
 \norm{\Delta_j^{-1/2}(U_B)_{j\cdot}}_2
       \le \frac{R}{n\sqrt p}\quad\text{whenever }x_j=1.
\end{equation}
Let $U$ contain orthonormal informative eigenvectors of $\widetilde A=WAW$, and put
\[
 H=G^{-1/2}UU^\top G^{1/2},\qquad
 H_B=\Delta^{-1/2}U_BU_B^\top\Delta^{1/2}.
\]
For every $\norm w_\infty\le C_w$, uniformly over the stated inputs,
\begin{equation}\label{restore:loo-main}
 (AHw)_i=x^\top H_Bw_{-i}+O(1).
\end{equation}
Implicit constants may depend on $L_D,R,C_w$ and the other fixed bounds
in the statement. No bound is assumed on unhit rows of $U_B$ or on $\max_j\Delta_j$.
\end{lemma}

\begin{proof}
For $D_i=0$ both row actions vanish. For $D_i>0$, put
$E_\Delta=\Delta^{-1/2}E\Delta^{-1/2}$, $h=\Delta^{-1/2}x$,
$\Pi_B^\perp=I-U_BU_B^\top$. Then $\norm{E_\Delta}\le C/(np)$ and
$\norm h^2\le CD_i/(np)$. The scalar-rescaled matrix $Q=G^{-1/2}AG^{-1/2}$ differs from
$\diag(0,Q_B)$ by $O({(np)}^{-1/2})$: its off-diagonal row has norm
at most $C\sqrt{D_i/(G_i{np})}$, and the lower renormalization has norm
$O({(np)}^{-1})$. Thus the full signal remains separated.

For an eigenvalue $\mu_r$ of $Q=\widetilde A/\tau_A$, write a full
generalized eigenvector as
\[
 G^{-1/2}u_r=\binom{t_r}{\Delta^{-1/2}(U_By_r+u_r^\perp)},
 \qquad U_B^\top u_r^\perp=0.
\]
The exact top and complementary equations are
\begin{align*}
 h^\top(U_By_r+u_r^\perp)&=\mu_r G_i t_r,\\
 (\mu_rI-N+\mu_r\Pi_B^\perp E_\Delta \Pi_B^\perp)u_r^\perp
   &=t_r\Pi_B^\perp h-\mu_r\Pi_B^\perp E_\Delta U_By_r.
\end{align*}
The inverse $T_r$ of the complementary operator is bounded. Therefore
\[
 (\mu_r G_i-(\Pi_B^\perp h)^\top T_r\Pi_B^\perp h)t_r
  =(U_B^\top h)^\top y_r
       -\mu_r(\Pi_B^\perp h)^\top T_r\Pi_B^\perp E_\Delta U_By_r.
\]
The denominator has magnitude at least $c{np}-CD_i/(np)\ge c'{np}$.
Here $G_i=\max(D_i,\tau_A)\asymp {np}$ is the diagonal weight in this formula.
The matrix $Y=(y_1,\ldots,y_K)$ and its inverse have bounded norm,
by the signal-space angle and $(I+E_\Delta)^{1/2}=I+O({(np)}^{-1})$.

The local row premise and the global threshold perturbation give
\[
 \norm{U_B^\top h}\le \frac{C_RD_i}{n\sqrt p},\qquad
 \norm{E_\Delta U_B}_F\le C_R\sqrt{D_i}/({np}\sqrt n)+C/n.
\]
The $C/n$ term is retained; it need not be supported on the neighbors.
Substitution yields, for $U_\perp=(u_1^\perp,\ldots,u_K^\perp)$,
\begin{equation}\label{restore:loo-refined}
 \norm t\le \frac{C_R}{n\sqrt p},\qquad
 \norm{U_\perp}\le C_R\{\sqrt{D_i}/({np}\sqrt n)+1/n\}.
\end{equation}
Let $V=\Delta^{-1/2}U_B$, $v_r=\Delta^{-1/2}u_r^\perp$ and
$f_r=Vy_r+v_r$. Write $t=(t_1,\ldots,t_K)^\top$, 
$v=(v_1,\ldots,v_K)$ and $f=(f_1,\ldots,f_K)$. Exact generalized orthonormality gives
\[
 Y^\top Y+U_\perp^\top U_\perp+G_i tt^\top+f^\top Ef=I.
\]
On the local part of $E$, use the hit-row bound; on its global part
use $\norm{E_\mathrm{glob}}\le Cp$ and
$\norm f\le C/\sqrt {np}$. Consequently
\begin{equation}\label{restore:loo-gram}
 \norm{YY^\top-I}\le C_R\{D_i/({n^2p})+1/n\}.
\end{equation}

For $w_V=V^\top\Delta w_{-i}$, the exact row expansion is
\[
 (AHw)_i=
 \sum_r(x^\top Vy_r+x^\top v_r)
 \{G_i t_rw_i+y_r^\top w_V+v_r^\top\Delta w_{-i}
                                      +f_r^\top Ew_{-i}\}.
\]
Only total degree is needed to bound global coefficient vectors. In
particular,
\begin{align*}
 \norm{x^\top V}&\le \frac{C_RD_i}{n\sqrt p},&
 \norm{w_V}&\le Cn\sqrt p,\\
 \norm{x^\top v}&\le C_R\{D_i/({(np)}^{3/2}\sqrt n)+\sqrt{D_i}/(n\sqrt {np})\},&
 \norm{v^\top\Delta w_{-i}}&\le C_R\{\sqrt{D_i/(np)}+\sqrt{p}\},\\
 \norm{G_i t}&\le C_R\sqrt{p},&
 \norm{f^\top Ew_{-i}}&\le C_R\{\frac{D_i}{n\sqrt p}+\sqrt{p}\}.
\end{align*}
For example the global part of the last bound follows from
$|f_r^\top E_\mathrm{glob}w_{-i}|
\le (Cp)C_w\sqrt n\norm{f_r}\le C\sqrt{p}$.
The leading term is $x^\top VYY^\top w_V$; its error relative to
$x^\top VV^\top\Delta w_{-i}$ is
$O(D_i^2/({n^2p})+D_i/n)$. Every other product is bounded by a fixed
multiple of
\[
 D_i/(np)+\sqrt{D_i/n}+D_i^2/({n^2p})+D_i/n+D_i\sqrt{D_i}/({np}\sqrt n).
\]
These quantities are bounded for $D_i\le L_D{np}$ and ${np}\le Cn$.
This proves the scalar comparison.
\end{proof}

For the empirical cap threshold $\tau=b_0\bar D$, the displayed decomposition of
$E$ is automatic on the two-sided total-degree event. Indeed, if
$\bar D_B$ is the mean degree of $B$, then
\[
 \tau_A-\tau_B=b_0(2D_i-\bar D_B)/n.
\]
First change each deleted degree by $x_j$ while keeping $\tau_B$;
this yields $E_\mathrm{loc}\in[0,\diag(x)]$. Then change the threshold;
this yields $\norm{E_\mathrm{glob}}\le C_{L_D}p$.

\begin{corollary}[Exact-degree reconstruction and simple hollowing]
\label{restore:loo-corrected}
Let $\widehat A^{\mathrm{c}}$ be the exact-degree hollow reconstruction in
\eqref{restore:reconstruction}. Then $\widehat A^{\mathrm{c}}\one=A\one$ and
$\diag(\widehat A^{\mathrm{c}})=0$ exactly.
For $\bar w_{-i}=(n-1)^{-1}\sum_{j\ne i}w_j$, the same local premises
give
\begin{equation}\label{restore:loo-corrected-action}
 (\widehat A^{\mathrm{c}} w)_i
 =x^\top\{H_B(w_{-i}-\bar w_{-i}\one)+\bar w_{-i}\one\}
       +O(1).
\end{equation}
\end{corollary}

\begin{proof}
The definition gives the exact row identity
\[
 (\widehat A^{\mathrm{c}} w)_i
 =[\widehat A(w-\bar w_{-i}\one)]_i+D_i\bar w_{-i}
       -\widehat A_{ii}(w_i-\bar w_{-i}).
\]
Apply the preceding lemma to $w-\bar w_{-i}\one$, whose entries
are bounded by $2C_w$. Equation~\eqref{restore:loo-refined} gives
$\norm{U_{i\cdot}}_2^2=G_i\norm t^2\le C_R/n$, so
$|\widehat A_{ii}|\le C_Rp$. The final term is therefore $O(p)=O(1)$,
proving \eqref{restore:loo-corrected-action}. The degree and diagonal
identities follow directly from \eqref{restore:reconstruction}.
\end{proof}

The independent coefficient vector in
\eqref{restore:loo-corrected-action} has perturbation
$(H_B-I)(w_{-i}-\bar w_{-i}\one)$. For block-constant contrasts it
therefore admits the tilted cancellation in
\cref{lem:sbm-floor-cancellation}.
For clarity, the hit-row exception is paid directly. On the
training event, $\Delta_j\ge c {np}$, so
\[
 \{j:\norm{\Delta_j^{-1/2}(U_B)_{j\cdot}}_2
                          >\frac{R}{n\sqrt p}\}
 \subseteq
 \{j:\sqrt n\norm{(U_B)_{j\cdot}}_2>R\sqrt c\}.
\]
\Cref{restore:tail-energy-lemma} bounds the expected cardinality of
this set on the training event by $Cn e^{-c'{np}\log R}$.
Conditioning only on the deleted graph and using $P_{z_i z_j}\le p$
therefore bounds the probability that the tested row hits it by
$C{np} e^{-c'{np}\log R}$. Choosing a sufficiently large fixed $R$
gives any prescribed fixed exponential rate $s$.
The upper-degree exception $D_i>L_D{np}$ has arbitrarily large exponential
rate after increasing the fixed $L_D$. Neither bound uses a union over
all possible tested targets.

Uniformity in $w$ is an algebraic stability statement. It does not make
full-graph fitted coefficients independent of $x$. The fitted-score
bound first replaces them by deterministic true contrasts; only then
is the independent sharp-tail bound applied.

\section{Block fitting and misclassification control}
\label{restore:fit}

Write $\Pi=UU^\top$. For the
proof, rescale the encoded adjacency matrix by its common threshold:
$Q=\widetilde A/\tau=G^{-1/2}AG^{-1/2}$. This changes the eigenvalue scale,
not the retained eigenspace. We work on the aggregate event
\begin{align}
 \norm{Q-G^{-1/2}{ZPZ^\top} G^{-1/2}}_{\op}
   &\le C/\sqrt {np},& \sum_iG_i&\le C{n^2p},
 \label{restore:fit-geometry-one}\\
 \sum_i\frac{(G_i-\tau)^2+(D_i-\sum_jP_{z_i z_j})^2}{G_i}
   &\le Cn,&
 \norm{(I-\Pi)G^{1/2}Z}_F&\le C\sqrt n,
 \label{restore:fit-geometry-two}\\
 \max_{k,\ell}\left|[Z^\top(A-{ZPZ^\top})Z]_{k\ell}\right|
   &\le Cn.&&\label{restore:fit-block-event}
\end{align}
We also have $G_i\ge\tau\asymp {np}$, $\sum_jP_{z_i z_j}\asymp {np}$,
$\bar D\asymp {np}$, and $\tau\ge\max_i\sum_jP_{z_i z_j}$. The first two
lines follow from \cref{prop:sparse-floor}; the block-total event is
paid in the risk proof. All bounds in this section are deterministic
on this event. Constants depend only on the fixed structural and
aggregate-event bounds and, for local estimates, on the fixed thresholds
$L_D,L_U$. They do not depend on the fitted partition. As in the
parameter-estimation and classification arguments of \citet{zhouli}, we display the dependence on misclassification rather
than absorbing it into a separately named error-radius constant.
Align the current labels once with the true labels. Then the incorrectly
assigned set $E=\{i:\widehat z_i\ne z_i\}$ has
$|E|=n\Mis(\widehat z,z)$. The iteration argument below preserves this
same label matching.
For current labels $\widehat z$, let
$\widehat Z_{ik}=\ind{\widehat z_i=k}$, and use
$n_k(\widehat z)$ for their community sizes.

An aligned population basis is
\[
 V=G^{-1/2}Z(Z^\top G^{-1}Z)^{-1/2}O,
\]
where $O$ is orthogonal. The signal gap gives
\begin{equation}\label{restore:fit-basis}
 \norm{U-V}_F\le C/\sqrt {np},\qquad
 \sqrt{G_i}\norm{V_{i\cdot}}_2\le C\sqrt{p}.
\end{equation}
Moreover $\norm Q_{\op}\le1$, since
$|x^\top Ax|\le\sum_iD_ix_i^2\le\sum_iG_ix_i^2$.
In particular $\norm{\Lambda/\tau}_{\op}\le1$, and
$\widehat A=G^{1/2}Q_KG^{1/2}$ for $Q_K=U(\Lambda/\tau)U^\top$.

For the row bounds in this section, put
$a_i=\sqrt{G_i}\norm{U_{i\cdot}}_2$.

\begin{lemma}[Weighted mass of incorrectly assigned vertices]
\label{restore:fit-wrong-mass}
For every set $E\subseteq[n]$,
\begin{equation}\label{restore:fit-wrong-mass-bounds}
 \sum_{i\in E}G_i\le C(n+np|E|),\qquad
 \sum_{i\in E}a_i\le C\left(|E|\sqrt p+\frac1{\sqrt p}\right).
\end{equation}
Also $\sum_i a_i\le Cn\sqrt p$. These bounds hold simultaneously
for all sets; in particular, the error set has
$|E|=n\Mis(\widehat z,z)$ under the common label matching.
\end{lemma}

\begin{proof}
The inequality $G_i\le2\tau+(G_i-\tau)^2/G_i$ and the aggregate
energy bound give $\sum_EG_i\le C(np|E|+n)$. By
\eqref{restore:fit-basis} and Cauchy--Schwarz,
\[
 \sum_{i\in E}a_i
 \le C|E|\sqrt p+
      \left(\sum_{i\in E}G_i\right)^{1/2}\norm{U-V}_F
 \le C\left(|E|\sqrt p+\sqrt{|E|+p^{-1}}\right).
\]
The inequality $\sqrt{p|E|+1}\le p|E|+1$ yields the second bound.
The global bound follows from
$\sum_i a_i\le(\sum_iG_i)^{1/2}\norm U_F$ and fixed $K$.
\end{proof}

\begin{lemma}[Radius-count initialization]
\label{restore:fit-initializer}
The initializer in \cref{restore:algorithm} satisfies
\[
 \Mis(\widehat z,z)\le \frac{C}{np}
\]
on the aggregate event. The comparison uses one common label permutation.
\end{lemma}

\begin{proof}
The normalized rows of $V$ are the $K$ rows of $O$, with pairwise
distance $\sqrt2$. Choose a fixed $B_0$ larger than twice the
uniform upper bound on $\tau/(np)$. Each row with $G_i>B_0{np}$
contributes at least $B_0{np}/4$ to $\sum_i(G_i-\tau)^2/G_i$, so
there are at most $\frac{C}{p}$ such rows. For every remaining row,
$\norm{V_{i\cdot}}_2\ge c/\sqrt n$. The normalization inequality
\[
 \left\|\frac u{\norm u_2}-\frac v{\norm v_2}\right\|_2
 \le\frac{2\norm{u-v}_2}{\norm v_2}
\]
(with any fixed unit direction when $u=0$) and
\eqref{restore:fit-basis} imply that the normalized empirical rows
$x_i$ have total squared distance at most $\frac{C}{p}$ from their true
centers $O_{z_i\cdot}$: the remaining rows satisfy this by the displayed
inequality, and the $O(\frac{1}{p})$ high-degree rows each contribute at most four.

Call a row good if $\norm{x_i-O_{z_i\cdot}}_2\le1/32$.
At most $m=\frac{C}{p}$ rows are bad, and $m<n/(3\beta K)$ for large $n$.
Inductively, all good rows of a selected block have been removed and
all good rows of unselected blocks remain. A radius-$1/8$ ball
centered at a good row of an unselected block contains every good row
of that block, hence more than $m$ remaining rows. The maximizing
ball must therefore contain a good row from an unselected block;
its representative is within $1/8+1/32=5/32$ of that center.
Every good row of that block lies within $3/16<1/4$ of the
representative. Good rows of any other block are farther than
$\sqrt2-3/16>1/4$, so deletion preserves the induction.
After $K$ selections, each good row is closer to its own
representative than to all others. Only the $m$ bad rows can be
misclassified. Dividing by $n$ gives the stated misclassification rate.
\end{proof}

For the remaining reconstruction estimates, write
$v=A\one-\widehat A\one$ and $R_1=\widehat A+v\one^\top/n$.
These are proof quantities computed from the same fixed reconstruction;
they are not additional algorithmic updates.

\begin{lemma}[Reconstructed block totals]
\label{restore:fit-block-totals}
The hollow, degree-preserving reconstruction in
\eqref{restore:reconstruction} satisfies
\begin{equation}\label{restore:fit-global-absolute}
 \sum_{i,j}|\widehat A^{\mathrm{c}}_{ij}|\le Cn^2p.
\end{equation}
For every estimated partition, under the common label matching,
\begin{equation}\label{restore:fit-total-accuracy}
 \left|[\widehat Z^\top\widehat A^{\mathrm{c}}\widehat Z]_{k\ell}
                  -n_kn_\ell P_{k\ell}\right|
 \le Cn\{1+np\Mis(\widehat z,z)\}.
\end{equation}
The omitted diagonal is included in the bound, and the constant is
uniform over the estimated partition.
\end{lemma}

\begin{proof}
With $r_i=(R_1)_{ii}$, the reconstruction formula is equivalently
\[
 \widehat A^{\mathrm{c}}=R_1+\frac{r\one^\top}{n-1}
                   -\frac n{n-1}\diag(r),
\]
because its off-diagonal entries are $\widehat A_{ij}+(v_i+\widehat A_{ii})/(n-1)$.
Let $Q_{\rm tail}=(I-\Pi)Q(I-\Pi)$. Spectral invariance gives
\[
 v=G^{1/2}Q_{\rm tail}(I-\Pi)G^{1/2}\one.
\]
The bounds $\norm{Q_{\rm tail}}_{\op}\le C/\sqrt{np}$,
$\norm{(I-\Pi)G^{1/2}\one}_2\le C\sqrt n$ and
$\sum_iG_i\le Cn^2p$ imply $\sum_i|v_i|\le Cn$.
Also
\[
 \sum_i|\widehat A_{ii}|\le\sum_iG_i\Pi_{ii}
 \le2\tau K+\sum_i\frac{(G_i-\tau)^2}{G_i}\le Cn.
\]
Consequently $\sum_i|(R_1)_{ii}|\le Cn$; hollowing has total
absolute cost at most $2\sum_i|(R_1)_{ii}|$.
Since $|\widehat A_{ij}|\le a_ia_j$, \cref{restore:fit-wrong-mass} bounds
its total absolute mass by $Cn^2p$. The rank-one degree correction
has absolute mass $\sum_i|v_i|$, proving
\eqref{restore:fit-global-absolute}.

Let $E$ be the incorrectly assigned set. Each nonzero row of
$\widehat Z-Z$ has squared norm two, so the block-space residual and
\cref{restore:fit-wrong-mass} give
\[
 \norm{(I-\Pi)G^{1/2}\widehat Z}_F^2
 \le 2\norm{(I-\Pi)G^{1/2}Z}_F^2
       +2\norm{G^{1/2}(\widehat Z-Z)}_F^2
 \le C(n+np|E|).
\]
The exact identity
\[
 \widehat Z^\top(\widehat A-A)\widehat Z
 =-\{(I-\Pi)G^{1/2}\widehat Z\}^\top Q_{\rm tail}
       \{(I-\Pi)G^{1/2}\widehat Z\}
\]
therefore has norm at most $C(n+np|E|)/\sqrt{np}$.
Moving the incorrectly assigned vertices changes any raw block total
by at most $C\sum_E D_i\le C(n+np|E|)$.
Together with \eqref{restore:fit-block-event}, this bounds the error
for $\widehat A$ by $C(n+np|E|)$.

The degree correction changes a block total by
$[\widehat Z^\top v]_k n_\ell(\widehat z)/n$.
Using spectral invariance again gives
\[
 \norm{\widehat Z^\top v}_2
 \le \frac{C\sqrt{n(n+np|E|)}}{\sqrt{np}}
 \le C(n+np|E|)
\]
for all sufficiently large $np$. The global $O(n)$ hollowing cost
has already been proved. Finally, $|E|=n\Mis(\widehat z,z)$ gives
\eqref{restore:fit-total-accuracy}.
\end{proof}

For a tested vertex $i$, the local regularity condition is
\begin{equation}\label{restore:fit-local-row}
 D_i\le L_Dnp,\qquad
 a_i=\sqrt{G_i}\norm{U_{i\cdot}}_2\le L_U\sqrt p,
\end{equation}
where $L_D,L_U$ are fixed. It is imposed only at the tested target.

\begin{lemma}[Local absolute mass]
\label{restore:fit-local-mass}
On \eqref{restore:fit-local-row},
\begin{equation}\label{restore:fit-local-mass-bounds}
 \sum_j|\widehat A^{\mathrm{c}}_{ij}|\le Cnp,\qquad
 \sum_{j\in E}|\widehat A^{\mathrm{c}}_{ij}|\le C(1+p|E|)
       \quad\text{for every }E\subseteq[n].
\end{equation}
For the incorrectly assigned set the second bound is
$C\{1+np\Mis(\widehat z,z)\}$.
\end{lemma}

\begin{proof}
The inequality $|\widehat A_{ij}|\le a_ia_j$, the local bound on $a_i$ and
\cref{restore:fit-wrong-mass} prove both claims for $\widehat A$.
They also give $|v_i|\le D_i+\sum_j|\widehat A_{ij}|\le Cnp$.
The rank-one correction adds at most $Cnp$ to the full row mass and
$Cp|E|$ to its mass over $E$. Finally,
$|(R_1)_{ii}|\le a_i^2+|v_i|/n\le Cp$; hollowing changes either
bound by at most $2|(R_1)_{ii}|$. This proves the assertion.
\end{proof}

\begin{proposition}[Parameter and score errors in terms of misclassification]
\label{restore:fit-ordinary}
\label{restore:fit-score}
There is a fixed $c>0$ such that, uniformly over current labels with
$\Mis(\widehat z,z)\le c$, the symmetric Bernoulli block fit is unique
and interior for all sufficiently large $n$, and
\begin{equation}\label{restore:fit-ordinary-parameter}
 \max_{k,\ell}|\widehat P_{k\ell}-P_{k\ell}|
 \le C\left\{\frac1n+p\Mis(\widehat z,z)\right\}.
\end{equation}
At every target satisfying \eqref{restore:fit-local-row}, the fitted
log-likelihood contrast obeys
\begin{equation}\label{restore:fit-ordinary-score}
 \begin{aligned}
 &\sum_{j\ne i}\left\{
 \widehat A^{\mathrm{c}}_{ij}\log\frac{\widehat P_{k\widehat z_j}}
                              {\widehat P_{\ell\widehat z_j}}
 +(1-\widehat A^{\mathrm{c}}_{ij})\log\frac{1-\widehat P_{k\widehat z_j}}
                                   {1-\widehat P_{\ell\widehat z_j}}
 \right\}\\
 &\qquad=[\widehat A^{\mathrm{c}} w_{k\ell}]_i
       +\sum_{j\ne i}\log\frac{1-P_{kz_j}}{1-P_{\ell z_j}}
       +O\{1+np\Mis(\widehat z,z)\}.
 \end{aligned}
\end{equation}
Here $w_{k\ell}$ is the true log-odds vector in
\eqref{eq:sbm-weights}. The constants are uniform over the current
labels. In particular, a misclassification rate of order $1/(np)$
gives parameter error $O(1/n)$ and a bounded score remainder.
\end{proposition}

\begin{proof}
For one unordered fitted block pair, write $E_{k\ell}$ for its ordered
reconstructed total and $N_{k\ell}$ for the number of ordered
 off-diagonal pairs. There are two directions when $k\ne\ell$ and
one ordered within-group total when $k=\ell$. Consequently
\cref{restore:fit-block-totals} and the exact pair counts give
\[
 \begin{aligned}
 E_{k\ell}
 &=(2-\ind{k=\ell})n_kn_\ell P_{k\ell}
        +O\bigl(n\{1+np\Mis(\widehat z,z)\}\bigr),\\
 N_{k\ell}
 &=(2-\ind{k=\ell})n_k(\widehat z)n_\ell(\widehat z)+O(n).
 \end{aligned}
\]
For the common label matching,
$|n_k(\widehat z)-n_k|\le n\Mis(\widehat z,z)$.
Choose the fixed $c$ small enough that every fitted group has size
at least $n/(2\beta K)$. Thus $N_{k\ell}\asymp n^2$, and
\[
 \left|\frac{E_{k\ell}}{N_{k\ell}}-P_{k\ell}\right|
 \le C\left\{\frac1n+p\Mis(\widehat z,z)\right\}.
\]
After decreasing $c$ if necessary, this error is a sufficiently small
fixed multiple of $p$ for all large $np$. The ratio is then positive,
comparable to $p$, and strictly inside the algorithm's clipping
interval. This uses the stipulated sufficiently small $c_0$ and
$p\to0$, a consequence of $np^2\le C_1$.
The criterion $E_{k\ell}\log q+(N_{k\ell}-E_{k\ell})\log(1-q)$
has the unique interior maximizer $q=E_{k\ell}/N_{k\ell}$,
proving \eqref{restore:fit-ordinary-parameter}.

On the positive comparable-probability class, the resulting log-odds
coefficient error is at most
$C\{(np)^{-1}+\Mis(\widehat z,z)\}$.
Multiplication by the absolute row mass $Cnp$ costs at most
$C\{1+np\Mis(\widehat z,z)\}$.
Changing the incorrectly assigned neighbor labels alters bounded
coefficients on $E$, at cost at most $C(1+p|E|)$ by
\cref{restore:fit-local-mass}.
For the nonedge intercept, bounded derivatives of $\log(1-q)$ give
parameter cost at most $C\{1+np\Mis(\widehat z,z)\}$;
wrong neighbor labels add at most $Cp|E|$.
As $|E|=n\Mis(\widehat z,z)$, these bounds prove
\eqref{restore:fit-ordinary-score}. They hold simultaneously over
all stated input labels; no independence of the fitted quantities
from the tested row is used.
\end{proof}

\begin{lemma}[Uniform misclassification control for finite hard refits]
\label{restore:fit-invariant-basin}
\label{restore:fit-invariant}
There are fixed constants $c>0$ and $C<\infty$ such that, on the
aggregate event and for all sufficiently large $n$, every input
satisfying $\Mis(\widehat z,z)\le c$ yields, after one block fit and
Bernoulli assignment,
\[
 \Mis\!\left(\mathcal L(\widehat A^{\mathrm{c}},\widehat P,\widehat z),z\right)
       \le\frac{C}{np}.
\]
Starting from the prescribed initializer, all iterates satisfy a
single bound $C/(np)$ under the same label matching.
For this initialized sequence, at every target satisfying
\eqref{restore:fit-local-row}, an error at any updated iterate implies
\begin{equation}\label{restore:fit-canonical-error}
 [\widehat A^{\mathrm{c}} w_{z_i\ell}]_i+
 \sum_{j\ne i}\log\frac{1-P_{z_i z_j}}{1-P_{\ell z_j}}\le C
 \quad\text{for some }\ell\ne z_i.
\end{equation}
The threshold is the same for all such iterates and for any finite,
possibly data-dependent stopping rule after at least one update.
\end{lemma}

\begin{proof}
Choose large fixed $L_D,R$, with $L_D$ larger than twice the uniform
upper bounds on $\sum_jP_{z_i z_j}/(np)$ and $\tau/(np)$.
The degree energy bounds $|\{i:D_i>L_Dnp\}|\le C/p$.
Also $\norm{V_{i\cdot}}_2\le C/\sqrt n$, so
\eqref{restore:fit-basis} gives
$|\{i:\sqrt n\norm{U_{i\cdot}}_2>R\}|\le C/p$.
Every other target satisfies \eqref{restore:fit-local-row} for a
fixed $L_U$, because $G_i\le Cnp$ there. These counting exceptions
are independent of the input labels.

Best rank-$K$ approximation gives
$\norm{Q_K-G^{-1/2}ZPZ^\top G^{-1/2}}_{\op}\le C/\sqrt{np}$,
whence
\[
 \norm{G^{-1/2}(\widehat A-ZPZ^\top)Z}_F^2
 \le \frac{C}{np}\norm{G^{1/2}Z}_F^2\le Cn.
\]
The earlier expression for $v$ implies
$\norm{G^{-1/2}v}_2^2\le C/p$.
Since $\norm{\one^\top Z/n}_2\le1$, the degree correction
preserves the bound:
\begin{equation}\label{restore:fit-basin-action}
 \norm{G^{-1/2}(R_1-ZPZ^\top)Z}_F^2\le Cn.
\end{equation}
Hollowing changes the action of a bounded block-constant vector
at row $i$ by at most $C|(R_1)_{ii}|$.
Because $\sum_i|(R_1)_{ii}|\le Cn$, it changes that action by more
than a prescribed small multiple of $np$ on at most $C/p$ rows.
This auxiliary set is used only for counting misclassifications;
it is not charged as a probability exception in the sharp risk proof.

For $k=z_i$, the true population contrast is
\[
 [ZPZ^\top w_{k\ell}]_i+
       \sum_{j\ne i}\log\frac{1-P_{kz_j}}{1-P_{\ell z_j}}.
\]
It is a sum of incident Bernoulli Kullback--Leibler divergences plus
an $O(p)$ completed-diagonal term. Row separation and community-size
bounds give a lower bound $c_{\rm gap}np$.
At a regular target, \cref{restore:fit-ordinary} bounds the fitted
remainder by $C\{1+np\Mis(\widehat z,z)\}$, with $C$ independent
of the input labels. First choose the fixed $c$ in the hypothesis
sufficiently small compared with this population margin, then take
$np$ sufficiently large. The remainder is at most $c_{\rm gap}np/4$.
Outside the auxiliary hollowing set, a new error therefore requires
\[
 \norm{[(R_1-ZPZ^\top)Z]_{i\cdot}}_2\ge c'np
\]
for a fixed $c'>0$. Summing the squared requirements divided by
$G_i\le Cnp$ and using \eqref{restore:fit-basin-action} bounds the
number of such errors by $C/p$. Including the counting exceptions
and dividing by $n$ proves the asserted output rate. Its constant
does not depend on the input misclassification rate within the
specified small interval.

The initializer has rate at most $C/(np)$ by
\cref{restore:fit-initializer}. Choose a single fixed upper constant
for this initializer bound and the one-update output bound, then
choose $n$ large enough that its ratio to $np$ is below the fixed
$c$ just chosen. Induction now controls every iterate without
introducing another radius parameter. The estimates use the same
true community numbering at every step, so the initial common
matching is retained.

Along this sequence $np\Mis(\widehat z,z)$ is uniformly bounded.
Equation~\eqref{restore:fit-ordinary-score} thus has a single bounded
remainder for all input iterates. A wrong assignment implies
\eqref{restore:fit-canonical-error}. The same event covers finite
adaptive stopping, with no probability union over iterations.
\end{proof}

\section{Exponential accuracy of the aggregate event}
\label{sec:high-range}

\begin{lemma}[Exponentially accurate aggregate geometry above a logarithmic scale]
\label{restore:high-aggregate}
Fix any $A_0>0$ and $s>0$. Under the fixed positive compact structural
conditions, if ${np}\ge A_0\log n$ and $np^2\le C_1$, there is an aggregate
event $\mathcal E_s$ with
\[
 \PP(\mathcal E_s^c)\le C_s e^{-s{np}}
\]
on which the capped-adjacency geometry, weighted degree-energy bounds,
two-sided total-degree bounds, and true block-total bounds used in
\cref{restore:fit-geometry-one,restore:fit-geometry-two,restore:fit-block-event}
hold for the full graph and every graph obtained by deleting at most three
vertices. The constants may depend on $A_0,s,C_1$ and the structural
bounds. No condition is imposed on the smallest raw degree or the
largest spectral row.
\end{lemma}

\begin{proof}
The centered bounded-entry norm inequality of \citet[Corollary 3.12 and Remark 3.13]{bvh} gives
\[
 \PP\{\norm{A-{ZPZ^\top}}_{\op}>L_s\sqrt {np}\}
 \le e^{-(s+2){np}}
\]
after choosing a sufficiently large fixed $L_s$. Indeed, the row
variance sum is $O({np})$, the entry bound is one, and the prefactor $n$
in the norm-tail inequality is absorbed using $\log n\le {np}/A_0$.
The completed-mean diagonal has norm $O(p)$ and is absorbed in the
same threshold. No minimum raw degree is imposed.

The Bernoulli transform of the total number of edges gives, with
failure probability at most $Ce^{-c{n^2p}}$,
\[
 c{n^2p}\le\sum_iD_i\le C{n^2p}.
\]
For each of the fixed $K^2$ true block totals, Bernstein's inequality
for the independent upper-triangular edges, whose coefficients have
absolute value at most two, gives
\[
 \PP\left\{\left|[Z^\top(A-{ZPZ^\top})Z]_{k\ell}\right|>Ln\right\}
 \le 2\exp\left\{-\frac{cL^2n}{{np}+L}\right\}.
\]
The deterministic completed diagonal costs $O({np})$ and is included
by increasing $L$. Since ${np}\to\infty$ and $np^2\le C_1$,
the exponent is at least $c'L^2{np}/C_1$ for all sufficiently large $n$.
Increasing the fixed $L=L_s$ therefore makes the failure of all block
totals at most $C_se^{-(s+2){np}}$. The total-degree failure is also
bounded by $Ce^{-(s+2){np}}$ for large $n$.

These three full-graph events imply the required events on every
fixed-depth minor without a degree-wise probability union. To see this,
let $J=[n]\setminus T$ with $|T|\le3$, set $A_J=A_{JJ}$ and
$Z_J=Z_{J\cdot}$, and recompute all degrees and the cap
threshold on $J$. The centered matrix $A_J-{Z_JPZ_J^\top}$ is a principal
submatrix of $A-{ZPZ^\top}$, so its norm is at most $L_s\sqrt {np}$.
Removing $T$ changes the total raw degree by at most $2|T|n$;
as ${np}\to\infty$, the two-sided bound by fixed multiples of ${n^2p}$
therefore persists uniformly over $T$. Removing $T$ changes a true
block total of $A-{ZPZ^\top}$ by $O(|T|n)$ deterministically, since
the entries have bounded absolute value. Its $O_s(n)$ bound persists
as well. The surviving group masses differ from their original
values by $O(1)$ and remain comparable to $n$.

The recomputed average degree is consequently comparable to ${np}$.
For the fixed sufficiently large cap multiplier $b_0$, its threshold
$\tau_J=b_0\overline D_J$ lies between positive fixed multiples of
${np}$ and above every completed row sum $\sum_{j\prime\in J}P_{z_j z_{j\prime}}$.
Thus $G_{J,j}=\max(D_{J,j},\tau_J)\ge c{np}$ and
$\sum_jG_{J,j}\le C{n^2p}$. The raw norm bound immediately gives
\begin{align*}
 \norm{G_J^{-1/2}(A_J-{Z_JPZ_J^\top})G_J^{-1/2}}_{\op}
   &\le C_s{(np)}^{-1/2},\\
 \sum_{j\in J}\frac{(D_{J,j}-\sum_{j\prime\in J}P_{z_j z_{j\prime}})^2}{G_{J,j}}
   &=\norm{G_J^{-1/2}(A_J-{Z_JPZ_J^\top})\one}_2^2
     \le C_sn.
\end{align*}
The cap exceeds the completed population degrees, so coordinatewise
\[
 \frac{(G_{J,j}-\tau_J)^2}{G_{J,j}}
 \le\frac{(D_{J,j}-\sum_{j\prime\in J}P_{z_j z_{j\prime}})^2}{G_{J,j}}.
\]
Indeed, the left side is zero when $D_{J,j}\le\tau_J$; otherwise
$0\le D_{J,j}-\tau_J\le D_{J,j}-\sum_{j\prime\in J}P_{z_j z_{j\prime}}$. This proves the
other weighted energy bound. The order-${np}$ signal gap of the
capped adjacency matrix (a fixed gap after dividing by $\tau_J$)
and the residual
$\norm{(I-U_JU_J^\top)G_J^{1/2}Z_J}_F\le C_s\sqrt n$
now follow by the deterministic Gram-matrix and projection arguments
in \cref{prop:sparse-floor}. Those arguments use only
$G_{J,j}\ge\tau_J\asymp {np}$, total degree, the displayed rescaled
noise bound, and the energy relative to $\tau_J$, all just proved.
Intersecting the stated aggregate consequences defines
$\mathcal E_s$. The three full-graph events imply it simultaneously
for every indicated minor, proving its claimed failure bound.

For a deleted graph used as independent training data, form the
analogous event using that graph and its further minors. The same
argument gives failure probability $C_se^{-s{np}}$, and this event is
measurable from the training graph alone. The full-graph common event
is used only for deterministic implications, never as a condition
that would preserve independence of an incident row.
\end{proof}


\section{Completion of the oracle-risk bound}
\label{sec:probability}

The final step uses a deleted graph only as a proof device. The
implemented procedure retains one full-graph eigensystem.
For each target $i$, put $B=A_{-i,-i}$ and recompute its degrees,
cap threshold, weights and projection $H_B$. Throughout this section, all constants
are uniform over the fixed structural class.

\begin{lemma}[Ordinary Bernoulli tilted cancellation]
\label{lem:sbm-floor-cancellation}
Let $w$ be a true block-constant Bernoulli log-odds vector on the
deleted graph and put
\[
 f=w-\bar w\one,\qquad \varepsilon=(H_B-I)f,
 \qquad \bar w=(n-1)^{-1}\one^\top w.
\]
On the deleted aggregate event, under the exact Chernoff tilt of
the raw-row oracle contrast, its Bernoulli probabilities $r_j$
satisfy
\begin{equation}\label{eq:sbm-floor-tilt}
 \left|\sum_{j\ne i}r_j\varepsilon_j\right|\le C,
 \qquad \sum_{j\ne i}r_j\varepsilon_j^2\le C.
\end{equation}
\end{lemma}

\begin{proof}
The tilted probabilities $r_j$ are constant within each true block
and satisfy $r_j\le Cp$. Write $\tau$ and $G$ for the threshold
and diagonal matrix recomputed from $B$. Since $r$ is block constant with $\norm r_\infty\le Cp$,
weighted orthogonality gives directly
\[
 r^\top\varepsilon
 =-\frac1\tau\big\langle(I-H_B)r,(I-H_B)f\big\rangle_G
       +\frac1\tau r^\top(\tau I-G)\varepsilon.
\]
The residual bounds give
$\norm{(I-H_B)r}_G\le Cp\sqrt n$ and
$\norm{(I-H_B)f}_G\le C\sqrt n$.
Also $\sum_j(G_j-\tau)^2/G_j\le Cn$, so
\[
 \norm{G^{-1/2}(\tau I-G)r}_2\le Cp\sqrt n,
 \qquad \norm\varepsilon_G\le C\sqrt n.
\]
Dividing either product by $\tau\asymp np$ gives $O(1)$.
This proves the first bound. Since $r_j\le Cp$ and
$\norm\varepsilon_2^2\le \frac{C}{p}$, the second bound follows as well.
The deterministic nonedge intercept remains in the oracle score
and is not absorbed into this cancellation.
When the cap is inactive on the deleted graph, $G=\tau I$ and
$H_B=\Pi_B=U_BU_B^\top$. The second term vanishes, leaving the
ordinary Euclidean adjacency identity
\[
 r^\top\varepsilon
 =-\big\langle(I-\Pi_B)r,(I-\Pi_B)f\big\rangle.
\]
\end{proof}

\begin{proof}[Proof of \cref{restore:sbm-theorem}]
Fix $A_0>0$. First consider ${np}\le A_0\log n$, with ${np}\to\infty$.
Let $\mathcal E$ collect the aggregate conditions
\eqref{restore:fit-geometry-one}--\eqref{restore:fit-block-event},
two-sided total degree, and their recomputed versions after deletion
of at most three vertices. For each prescribed fixed $r>0$,
\cref{prop:sparse-floor} gives these geometric conditions with
failure probability at most $C_rn^{-r}$, by increasing its
probability parameter to cover the polynomial number of minors.
For any true block total, independent-edge Bernstein gives
\[
 \PP\{|[Z^\top(A-{ZPZ^\top})Z]_{k\ell}|>Ln\}
 \le2\exp\{-cL^2n/({np}+L)\}.
\]
In the present range this is smaller than any fixed inverse power
of $n$, as is the total-degree failure. The same holds for the
required minors. For a target $i$, define $\mathcal E_i$ using
only its deleted graph and its further deletions. Then
\begin{equation}\label{eq:sbm-geometry-payment}
 \PP(\mathcal E^c)+\frac1n\sum_i\PP(\mathcal E_i^c)
 \le C_rn^{-r}.
\end{equation}

On $\mathcal E$, \cref{restore:fit-initializer,restore:fit-invariant}
give the initializer and every subsequent iterate the same bound
$\Mis(\widehat z,z)\le C/(np)$. At target
$i$ impose only $D_i\le L_D{np}$, its local spectral-row condition
\eqref{restore:fit-local-row}, and the requirement that the
incident row hits no exceptional deleted spectral row. By
\cref{restore:tail-energy-lemma,restore:tail-hit-corollary,restore:tail-target-corollary},
the average probability cost of these spectral exceptions is at
most $C_se^{-s{np}}$ for any prescribed fixed $s$, on choosing fixed
thresholds sufficiently large. Indeed, $G_j\ge\tau\asymp {np}$ gives
\eqref{restore:loo-hit} outside a deleted exceptional set
$\{j:\sqrt n\norm{(U_B)_{j\cdot}}_2>R\}$, whose expected hit
mass is controlled by \eqref{restore:tail-coefficient-count}.
The degree tail has the same bound by the single-row Bernoulli
Chernoff inequality. Polynomial factors in ${np}$ are absorbed by
first choosing a larger exponential rate. These are averages of
target probabilities, without a union over all tested vertices.

On the resulting local events, the deterministic fitted-score
bound and \cref{restore:loo-corrected} imply that an aligned error
at $i$ requires, for some $\ell\ne k=z_i$,
\begin{equation}\label{eq:sbm-independent-score}
 T_{k\ell}+\sum_{j\ne i}A_{ij}\varepsilon_j\le C,
 \qquad \varepsilon=(H_B-I)(w_{k\ell,-i}-\bar w_{k\ell,-i}\one).
\end{equation}
Here $T_{k\ell}$ is realized from this target's actual row, with its
full nonedge term and $n_k-1$ own-community trials. Its unconditional
law is the community-$k$ law defined in Section~\ref{sec:statements};
no independent substitute is introduced. The coefficients $\varepsilon$
are measurable from $B$ and independent of $(A_{ij})_{j\ne i}$. To apply the probability
bound, first use this deterministic inclusion, then drop the
full-graph event and condition only on $B$ on $\mathcal E_i$.
No independence is asserted conditional on $\mathcal E$.
The score threshold is fixed uniformly over all initialized refits,
because $1+np\Mis(\widehat z,z)$ is uniformly bounded by
\cref{restore:fit-invariant}.

The tilted mean and quadratic bounds in
\cref{lem:sbm-floor-cancellation} give the moment hypotheses of
\cref{lem:sparse-truncated-transfer}. Its exceptional-coefficient
hit term is at most $C_se^{-s{np}}$ by
\cref{restore:tail-hit-corollary}, after selecting a fixed
coefficient threshold. That truncation is used only in the proof.
Apply the sharp transfer to \eqref{eq:sbm-independent-score} using
the exact self-excluded row law. By the one-coordinate comparison in
Appendix~\ref{sec:sharp}, its exponential risk factor can be replaced
by $e^{-I_{k\ell}}$ at a uniformly bounded multiplicative cost.
Averaging targets and taking the fixed-$K$ union over competitors,
we obtain
\begin{equation}\label{eq:sbm-final-payment}
 \begin{split}
 \E\Mis(\widehat z,z)
 \lesssim{}&\sum_{k=1}^K\frac{n_k}{n}
                     \sum_{\ell\ne k}e^{-I_{k\ell}}\\
 &+e^{-snp}
       +\PP(\mathcal E^c)+\frac1n\sum_i\PP(\mathcal E_i^c).
 \end{split}
\end{equation}
Here the implicit factor may also depend on the chosen fixed
$s$ and the associated fixed local thresholds, never on $n$ or the
number of refits. We grouped the local testing terms by $k=z_i$;
the graph-dependent exceptional events retain their vertex indices.
The comparison in \eqref{eq:oracle-pairs} bounds the testing term by
a constant multiple of $R_{\mathrm{oracle}}$, for the specified tie rule.

On the compact structural class, choose a fixed $C$ such that
\[
 I_{k\ell}-\frac12\log(np)\le Cnp
\]
uniformly over the tested pairs. Consequently
\[
 R_{\mathrm{oracle}}\gtrsim\exp\{-Cnp-\tfrac12\log(np)\}.
\]
Choose $s>C+1$ and $r>A_0(C+1)+2$. When
$np\le A_0\log n$,
\[
 \sqrt{np}\left(e^{-(s-C)np}+n^{-r}e^{Cnp}\right)\longrightarrow0,
\]
which pays both remainders in \eqref{eq:sbm-final-payment}.

For ${np}>A_0\log n$, use \cref{restore:high-aggregate} in place of
the polynomially accurate geometry. It gives
$\PP(\mathcal E^c)+n^{-1}\sum_i\PP(\mathcal E_i^c)\le C_{s'}e^{-s'{np}}$ for any fixed
$s'>0$, with training events still measurable from the respective
deleted graphs. The row-tail proof uses deterministic aggregate
geometry and applies innovation concentration before restricting
to that event. Its estimates, the local LOO comparison and the
fitting argument therefore remain valid for this stronger event.
Equation \eqref{eq:sbm-final-payment} follows with the new geometry
remainder. Choosing $s,s'>C+1$ pays both exponential terms.
The two proof ranges cover every sequence ${np}\to\infty$ under
$np^2\le C_1$, proving the theorem without a logarithmic degree
lower bound.
\end{proof}

\bibliographystyle{plain}
\bibliography{references}
\end{document}

